\ifnum\EditionMode=1\relax
\clearpage
\chapter*{Apparat des corrections}
\addcontentsline{toc}{chapter}{Apparat des corrections}
Chaque numéro correspond à un lieu précis du texte corrigé. Les fragments qui ne peuvent pas être composés isolément (délimiteur ouvert, fragment de diagramme) sont reproduits exactement en notation TeX. Les réparations de transcription communes aux trois formes ne figurent pas ici.
\begingroup\small
\begin{description}
\item[\hypertarget{apparatus:FR-0001}{1}] Page \hyperlink{printpage:fr:III}{III}\quad\texttt{SGA45-CORR-0001}
\par\noindent\textbf{Imprimé:} digression
\par\noindent\textbf{Corrigé:} digressions
\par\noindent Plural series and certaines require digressions; print has singular.
\smallskip
\item[\hypertarget{apparatus:FR-0002}{2}] Page \hyperlink{printpage:fr:1}{1}\quad\texttt{SGA45-CORR-0002}
\par\noindent\textbf{Imprimé:} cohomologiques
\par\noindent\textbf{Corrigé:} cohomologique
\par\noindent Singular dimension requires a singular adjective.
\smallskip
\item[\hypertarget{apparatus:FR-0003}{3}] Page \hyperlink{printpage:fr:10}{10}\quad\texttt{SGA45-CORR-0004}
\par\noindent\textbf{Imprimé:} plat est surjectif
\par\noindent\textbf{Corrigé:} plat et surjectif
\par\noindent The printed est is a typographical substitution for conjunction et.
\smallskip
\item[\hypertarget{apparatus:FR-0004}{4}] Page \hyperlink{printpage:fr:10}{10}\quad\texttt{S45-SOL-0027}
\par\noindent\textbf{Imprimé:} Pour toute famille
\par\noindent\textbf{Corrigé:} Pour toute famille finie
\par\noindent The containment claim is valid for finite families; for infinite families it fails in an infinite free module.
\smallskip
\item[\hypertarget{apparatus:FR-0005}{5}] Page \hyperlink{printpage:fr:12}{12}\quad\texttt{S45-SOL-0033}
\par\noindent\textbf{Imprimé:} \(M_U=\varphi^* M'\)
\par\noindent\textbf{Corrigé:} \(M_V=\varphi^* M'\)
\par\noindent The pullback along V→U lives on V.
\smallskip
\item[\hypertarget{apparatus:FR-0006}{6}] Page \hyperlink{printpage:fr:13}{13}\quad\texttt{S45-SOL-0035}
\par\noindent\textbf{Imprimé:} \(\partial_0\partial_1\)
\par\noindent\textbf{Corrigé:} \(\partial_1\partial_0\)
\par\noindent The usual coface relation for right-to-left composition requires the reversed left indices.
\smallskip
\item[\hypertarget{apparatus:FR-0007}{7}] Page \hyperlink{printpage:fr:14}{14}\quad\texttt{S45-SOL-0038}
\par\noindent\textbf{Imprimé:} \(v'\in V\)
\par\noindent\textbf{Corrigé:} \(v'\in V'\)
\par\noindent The arbitrary basis vector is in V', not necessarily its invariant subspace V.
\smallskip
\item[\hypertarget{apparatus:FR-0008}{8}] Page \hyperlink{printpage:fr:16}{16}\quad\texttt{S45-SOL-0045}
\par\noindent\textbf{Imprimé:} telle que
\par\noindent\textbf{Corrigé:} tel que
\par\noindent Agreement with masculine tout s.
\smallskip
\item[\hypertarget{apparatus:FR-0009}{9}] Page \hyperlink{printpage:fr:17}{17}\quad\texttt{SGA45-CORR-0006}
\par\noindent\textbf{Imprimé:} 6.5
\par\noindent\textbf{Corrigé:} 6.6
\par\noindent Duplicated paragraph number; the next number is 6.6.
\smallskip
\item[\hypertarget{apparatus:FR-0010}{10}] Page \hyperlink{printpage:fr:20}{20}\quad\texttt{S45-SOL-0057}
\par\noindent\textbf{Imprimé:} \hyperref[I:1-6-6]{I.6.5}
\par\noindent\textbf{Corrigé:} \hyperref[I:1-6-6]{I.6.6}
\par\noindent Reference follows the corrected torsor paragraph number.
\smallskip
\item[\hypertarget{apparatus:FR-0011}{11}] Page \hyperlink{printpage:fr:22}{22}\quad\texttt{SGA45-CORR-0008}
\par\noindent\textbf{Imprimé:} induit
\par\noindent\textbf{Corrigé:} induits
\par\noindent Plural homomorphismes governs the participle.
\smallskip
\item[\hypertarget{apparatus:FR-0012}{12}] Page \hyperlink{printpage:fr:24}{24}\quad\texttt{S45-SOL-0085}
\par\noindent\textbf{Imprimé:} en $y$
\par\noindent\textbf{Corrigé:} en $\bar y$
\par\noindent Strict localization is at the chosen geometric point.
\smallskip
\item[\hypertarget{apparatus:FR-0013}{13}] Page \hyperlink{printpage:fr:25}{25}\quad\texttt{S45-SOL-0089}
\par\noindent\textbf{Imprimé:} polynome $F$
\par\noindent\textbf{Corrigé:} polynôme $F$
\par\noindent Circumflex missing in the printed spelling.
\smallskip
\item[\hypertarget{apparatus:FR-0014}{14}] Page \hyperlink{printpage:fr:28}{28}\quad\texttt{S45-SOL-0094}
\par\noindent\textbf{Imprimé:} \(A_K\)
\par\noindent\textbf{Corrigé:} \(A_{K'}\)
\par\noindent The indicated extension of scalars is to K′.
\smallskip
\item[\hypertarget{apparatus:FR-0015}{15}] Page \hyperlink{printpage:fr:30}{30}\quad\texttt{S45-SOL-0103}
\par\noindent\textbf{Imprimé:} \((\bar K\otimes_K K')^*/\bar k^*\)
\par\noindent\textbf{Corrigé:} \((\bar K\otimes_K K')^*/\bar K^*\)
\par\noindent The closure is that of K; lowercase k is undefined.
\smallskip
\item[\hypertarget{apparatus:FR-0016}{16}] Page \hyperlink{printpage:fr:32}{32}\quad\texttt{S45-SOL-0114}
\par\noindent\textbf{Imprimé:} \(g(X_{i j})\)
\par\noindent\textbf{Corrigé:} \(g(U_{i j})\)
\par\noindent The variables introduced in the previous sentence are U\_\{ij\}.
\smallskip
\item[\hypertarget{apparatus:FR-0017}{17}] Page \hyperlink{printpage:fr:32}{32}\quad\texttt{SGA45-CORR-0011B}
\par\noindent\textbf{Imprimé:} \(N_{K/k}\)
\par\noindent\textbf{Corrigé:} \(N_{K/k(X)}\)
\par\noindent Only K/k(X) is the indicated finite extension.
\smallskip
\item[\hypertarget{apparatus:FR-0018}{18}] Page \hyperlink{printpage:fr:45}{45}\quad\texttt{S45-SOL-0146}
\par\noindent\textbf{Imprimé:} \(\prod_{p_{i*}} C_i\)
\par\noindent\textbf{Corrigé:} \(\prod p_{i*} C_i\)
\par\noindent The product is of the direct images p\_\{i*\}C\_i, as in (3.2)(ii).
\smallskip
\item[\hypertarget{apparatus:FR-0019}{19}] Page \hyperlink{printpage:fr:45}{45}\quad\texttt{S45-ASTRA-A-0021}
\par\noindent\textbf{Imprimé:} $n$ un entier $\geqslant 0$
\par\noindent\textbf{Corrigé:} $n$ un entier $>0$
\par\noindent The reduction uses torsion coefficients Z/n and then reduces n to a prime power; n=0 would give Z and is outside that argument.
\smallskip
\item[\hypertarget{apparatus:FR-0020}{20}] Page \hyperlink{printpage:fr:49}{49}\quad\texttt{SGA45-CORR-0015}
\par\noindent\textbf{Imprimé:} le 
\par\noindent\textbf{Corrigé:} \emph{\textnormal{[omission]}}
\par\noindent The preceding page already ends with le; remove only the duplicated article in corrected/reading.
\smallskip
\item[\hypertarget{apparatus:FR-0021}{21}] Page \hyperlink{printpage:fr:49}{49}\quad\texttt{S45-SOL-0154}
\par\noindent\textbf{Imprimé:} \(\eR^q f_!\)
\par\noindent\textbf{Corrigé:} \(\eR^q f_! \sF\)
\par\noindent The displayed equality defines the value at F.
\smallskip
\item[\hypertarget{apparatus:FR-0022}{22}] Page \hyperlink{printpage:fr:49}{49}\quad\texttt{SGA45-CORR-0014}
\par\noindent\textbf{Imprimé:} \(\eR^q f_*(j_! \sF)\)
\par\noindent\textbf{Corrigé:} \(\eR^q \bar f_*(j_! \sF)\)
\par\noindent The argument is a sheaf on the compactification, whose map is barred f.
\smallskip
\item[\hypertarget{apparatus:FR-0023}{23}] Page \hyperlink{printpage:fr:54}{54}\quad\texttt{SGA45-CORR-0019}
\par\noindent\textbf{Imprimé:} \(\h^0\left(X_t^x,\dZ/n\right)\)
\par\noindent\textbf{Corrigé:} \(\h^0\left(\widetilde X_t^x,\dZ/n\right)\)
\par\noindent The expression concerns the vanishing-cycle variety already defined with a tilde.
\smallskip
\item[\hypertarget{apparatus:FR-0024}{24}] Page \hyperlink{printpage:fr:56}{56}\quad\texttt{S45-ASTRA-A-0020}
\par\noindent\textbf{Imprimé:} \(\h^*(X,\varepsilon_*' \dZ/n)\)
\par\noindent\textbf{Corrigé:} \(\h^*(\widetilde X,\varepsilon_*' \dZ/n)\)
\par\noindent The direct image by epsilon-prime is a sheaf on X-tilde in the preceding Cartesian diagram.
\smallskip
\item[\hypertarget{apparatus:FR-0025}{25}] Page \hyperlink{printpage:fr:59}{59}\quad\texttt{SGA45-CORR-0021}
\par\noindent\textbf{Imprimé:} algèbres locales étale
\par\noindent\textbf{Corrigé:} algèbres locales étales
\par\noindent Plural agreement with algèbres locales.
\smallskip
\item[\hypertarget{apparatus:FR-0026}{26}] Page \hyperlink{printpage:fr:61}{61}\quad\texttt{S45-SOL-0209}
\par\noindent\textbf{Imprimé:} \(\widehat X=A\llbracket T\rrbracket\)
\par\noindent\textbf{Corrigé:} \(\widehat X=\spec A\llbracket T\rrbracket\)
\par\noindent The completed base is a scheme; Spec is missing from the printed equality.
\smallskip
\item[\hypertarget{apparatus:FR-0027}{27}] Page \hyperlink{printpage:fr:64}{64}\quad\texttt{S45-SOL-0217}
\par\noindent\textbf{Imprimé:} \(\dA_S^{n-1}\)
\par\noindent\textbf{Corrigé:} \(\dA_S^{N-1}\)
\par\noindent N is the relative dimension; n is the torsion modulus.
\smallskip
\item[\hypertarget{apparatus:FR-0028}{28}] Page \hyperlink{printpage:fr:70}{70}\quad\texttt{S45-SOL-0244}
\par\noindent\textbf{Imprimé:} corps de classe géométrique
\par\noindent\textbf{Corrigé:} corps de classes géométrique
\par\noindent The established French mathematical term is corps de classes; the final s is missing.
\smallskip
\item[\hypertarget{apparatus:FR-0029}{29}] Page \hyperlink{printpage:fr:70}{70}\quad\texttt{S45-SOL-0248}
\par\noindent\textbf{Imprimé:} Soient $X$ une variété
\par\noindent\textbf{Corrigé:} Soit $X$ une variété
\par\noindent Only one object X is introduced, so the introductory verb must be singular.
\smallskip
\item[\hypertarget{apparatus:FR-0030}{30}] Page \hyperlink{printpage:fr:72}{72}\quad\texttt{S45-SOL-0257}
\par\noindent\textbf{Imprimé:} du paragraphe précédents
\par\noindent\textbf{Corrigé:} du paragraphe précédent
\par\noindent The adjective agrees with the singular noun paragraphe.
\smallskip
\item[\hypertarget{apparatus:FR-0031}{31}] Page \hyperlink{printpage:fr:74}{74}\quad\texttt{S45-ASTRA-A-0029}
\par\noindent\textbf{Imprimé:} 1959
\par\noindent\textbf{Corrigé:} 1969
\par\noindent The original journal record gives volume36 (1969), DOI10.1007/BF02684596; the printed bibliography has the wrong decade. Verified at https://www.numdam.org/item/PMIHES\_1969\_\_36\_\_23\_0/.
\smallskip
\item[\hypertarget{apparatus:FR-0032}{32}] Page \hyperlink{printpage:fr:74}{74}\quad\texttt{S45-ASTRA-A-0030}
\par\noindent\textbf{Imprimé:} 25--58
\par\noindent\textbf{Corrigé:} 23--58
\par\noindent The journal record and the article first page give23–58, not25–58. DOI10.1007/BF02684596; https://www.numdam.org/item/PMIHES\_1969\_\_36\_\_23\_0/.
\smallskip
\item[\hypertarget{apparatus:FR-0033}{33}] Page \hyperlink{printpage:fr:74}{74}\quad\texttt{S45-ASTRA-A-0033}
\par\noindent\textbf{Imprimé:} 266
\par\noindent\textbf{Corrigé:} 256
\par\noindent The same work is correctly numbered256 in Ariane p3; its primary Bourbaki record confirms256: https://www.numdam.org/item/SB\_1962-1964\_\_8\_\_189\_0/.
\smallskip
\item[\hypertarget{apparatus:FR-0034}{34}] Page \hyperlink{printpage:fr:75}{75}\quad\texttt{S45-ASTRA-A-0031}
\par\noindent\textbf{Imprimé:} 1--12
\par\noindent\textbf{Corrigé:} 1--42
\par\noindent The original journal record gives Annales de l’Institut Fourier6 (1956),1–42, DOI10.5802/aif.59; https://www.numdam.org/item/?id=AIF\_1956\_\_6\_\_1\_0. The printed final page12 is wrong.
\smallskip
\item[\hypertarget{apparatus:FR-0035}{35}] Page \hyperlink{printpage:fr:80}{80}\quad\texttt{S45-SOL-0283}
\par\noindent\textbf{Imprimé:} Ceci à
\par\noindent\textbf{Corrigé:} Ceci a
\par\noindent The visible grave accent gives the preposition à, but the subject Ceci requires the verb a; French spelling correction has no separate English variant.
\smallskip
\item[\hypertarget{apparatus:FR-0036}{36}] Page \hyperlink{printpage:fr:82}{82}\quad\texttt{S45-SOL-0286}
\par\noindent\textbf{Imprimé:} modules constructibles
\par\noindent\textbf{Corrigé:} modules constructible
\par\noindent Constructibility is a property of the singular faisceau, not of its coefficient modules; remove the printed agreement error only in corrected forms.
\smallskip
\item[\hypertarget{apparatus:FR-0037}{37}] Page \hyperlink{printpage:fr:84}{84}\quad\texttt{S45-SOL-0305}
\par\noindent\textbf{Imprimé:} est la $\dQ_\ell$-espace
\par\noindent\textbf{Corrigé:} est le $\dQ_\ell$-espace
\par\noindent The print visibly uses la, but espace is masculine; no semantic English variant.
\smallskip
\item[\hypertarget{apparatus:FR-0038}{38}] Page \hyperlink{printpage:fr:85}{85}\quad\texttt{S45-SOL-0309}
\par\noindent\textbf{Imprimé:} $\dZ_\ell$-faisceau sur $Y$, alors
\par\noindent\textbf{Corrigé:} $\dZ_\ell$-faisceau sur $X$, alors
\par\noindent For the declared f:X→Y, the input to R\textasciicircum{}q f\_! must be a sheaf on X; the print visibly has Y in that hypothesis. The output remains on Y.
\smallskip
\item[\hypertarget{apparatus:FR-0039}{39}] Page \hyperlink{printpage:fr:90}{90}\quad\texttt{S45-SOL-0326}
\par\noindent\textbf{Imprimé:} des sous-catégorie
\par\noindent\textbf{Corrigé:} des sous-catégories
\par\noindent The plural article des and three following derived subcategories require the plural noun; print visibly omits final s.
\smallskip
\item[\hypertarget{apparatus:FR-0040}{40}] Page \hyperlink{printpage:fr:92}{92}\quad\texttt{S45-SOL-0342}
\par\noindent\textbf{Imprimé:} $\gr_F^p(\Lambda)$ soient
\par\noindent\textbf{Corrigé:} $\gr_F^p(K)$ soient
\par\noindent The source visibly prints Lambda in the graded object. The surrounding pair (K,F), preceding graded terms, and essential-image criterion require the graded pieces of K.
\smallskip
\item[\hypertarget{apparatus:FR-0041}{41}] Page \hyperlink{printpage:fr:93}{93}\quad\texttt{S45-SOL-0345}
\par\noindent\textbf{Imprimé:} \(\tor_k(N,K^r/\im d)\)
\par\noindent\textbf{Corrigé:} \(\tor_k(N,K^{-r}/\im d)\)
\par\noindent The visible exponent r lacks a minus sign; the preceding free resolution starts in degree -r and computes Tor with its cokernel.
\smallskip
\item[\hypertarget{apparatus:FR-0042}{42}] Page \hyperlink{printpage:fr:93}{93}\quad\texttt{S45-SOL-0346}
\par\noindent\textbf{Imprimé:} le module $K^r/\im d$
\par\noindent\textbf{Corrigé:} le module $K^{-r}/\im d$
\par\noindent This is the cokernel of the resolution in degree -r; the scan repeats the missing minus sign.
\smallskip
\item[\hypertarget{apparatus:FR-0043}{43}] Page \hyperlink{printpage:fr:93}{93}\quad\texttt{S45-SOL-0347}
\par\noindent\textbf{Imprimé:} \(0 \to K^r/\im(d) \to K^{r+1} \to K^{r+2}\)
\par\noindent\textbf{Corrigé:} \(0 \to K^{-r}/\im(d) \to K^{-r+1} \to K^{-r+2}\)
\par\noindent All three displayed exponents visibly use positive r; the truncation must start at -r, consistently with the resolution and Tor-dimension bound.
\smallskip
\item[\hypertarget{apparatus:FR-0044}{44}] Page \hyperlink{printpage:fr:94}{94}\quad\texttt{S45-SOL-0351}
\par\noindent\textbf{Imprimé:} Si un complexe de
\par\noindent\textbf{Corrigé:} Si un complexe $\sK$ de
\par\noindent The print omits the name of the complex but immediately uses K in H\textasciicircum{}i(K) and K′→K; supply the missing name in corrected forms.
\smallskip
\item[\hypertarget{apparatus:FR-0045}{45}] Page \hyperlink{printpage:fr:94}{94}\quad\texttt{S45-SOL-0353}
\par\noindent\textbf{Imprimé:} construire $\sK^n$ plat et constructible et $v:\sK^n\to \sA$
\par\noindent\textbf{Corrigé:} construire ${\sK'}^n$ plat et constructible et $v:{\sK'}^n\to \sA$
\par\noindent The scan twice gives unprimed K\textasciicircum{}n, although the induction constructs the terms of K′.
\smallskip
\item[\hypertarget{apparatus:FR-0046}{46}] Page \hyperlink{printpage:fr:95}{95}\quad\texttt{S45-SOL-0355}
\par\noindent\textbf{Imprimé:} Soit $F:X\to Y$ un morphisme séparé
\par\noindent\textbf{Corrigé:} Soit $f:X\to Y$ un morphisme séparé
\par\noindent Capital F in the theorem is inconsistent with the subsequent functor f; the diagram itself correctly has lowercase barred f and is not changed.
\smallskip
\item[\hypertarget{apparatus:FR-0047}{47}] Page \hyperlink{printpage:fr:95}{95}\quad\texttt{S45-SOL-0357}
\par\noindent\textbf{Imprimé:} $\eR f_!\sK\in \ob \eD^b(X,\Lambda)$
\par\noindent\textbf{Corrigé:} $\eR f_!\sK\in \ob \eD^b(Y,\Lambda)$
\par\noindent The scan prints X, but pushforward under f:X→Y is a complex on Y, as the theorem states.
\smallskip
\item[\hypertarget{apparatus:FR-0048}{48}] Page \hyperlink{printpage:fr:95}{95}\quad\texttt{S45-SOL-0359}
\par\noindent\textbf{Imprimé:} $\eR f_! = \eR f_*\circ j_!$
\par\noindent\textbf{Corrigé:} $\eR f_! = \eR \bar f_*\circ j_!$
\par\noindent The compactification uses proper map bar f. The bottom line drops the bar, unlike the definition above; corrected forms restore it.
\smallskip
\item[\hypertarget{apparatus:FR-0049}{49}] Page \hyperlink{printpage:fr:96}{96}\quad\texttt{S45-SOL-0361}
\par\noindent\textbf{Imprimé:} Si $L$ un
\par\noindent\textbf{Corrigé:} Si $L$ est un
\par\noindent The conditional clause lacks its verb est in the print.
\smallskip
\item[\hypertarget{apparatus:FR-0050}{50}] Page \hyperlink{printpage:fr:96}{96}\quad\texttt{S45-SOL-0362}
\par\noindent\textbf{Imprimé:} $\Lambda$ libre de type fini
\par\noindent\textbf{Corrigé:} $L$ libre de type fini
\par\noindent The free right module being varied is L; Lambda is the coefficient ring. The following passage-to-the-limit argument concerns L.
\smallskip
\item[\hypertarget{apparatus:FR-0051}{51}] Page \hyperlink{printpage:fr:97}{97}\quad\texttt{S45-SOL-0370}
\par\noindent\textbf{Imprimé:} \(\sK\in\ob_{\text{ctf}}(X,\Lambda)\)
\par\noindent\textbf{Corrigé:} \(\sK\in\ob\eD_\text{ctf}(X,\Lambda)\)
\par\noindent The scan omits the category letter D after Ob; the object must belong to the derived category with constructible finite-Tor-dimension cohomology.
\smallskip
\item[\hypertarget{apparatus:FR-0052}{52}] Page \hyperlink{printpage:fr:98}{98}\quad\texttt{S45-SOL-0378}
\par\noindent\textbf{Imprimé:} pour 
$\dZ_\ell$-faisceau
\par\noindent\textbf{Corrigé:} pour un 
$\dZ_\ell$-faisceau
\par\noindent The print omits the indefinite article before the singular sheaf.
\smallskip
\item[\hypertarget{apparatus:FR-0053}{53}] Page \hyperlink{printpage:fr:98}{98}\quad\texttt{S45-SOL-0379}
\par\noindent\textbf{Imprimé:} $K_n=\eR\Gamma_c(X,\sF\otimes\dZ/\ell^n)$
\par\noindent\textbf{Corrigé:} $K_n=\eR\Gamma_c(X,\sF\otimes\dZ/\ell^{n+1})$
\par\noindent The coefficient exponent n contradicts the following category D\_parf(Z/l\textasciicircum{}\{n+1\}), formula4.11.1, and trace computation on p.99; restore n+1 only in corrected forms.
\smallskip
\item[\hypertarget{apparatus:FR-0054}{54}] Page \hyperlink{printpage:fr:99}{99}\quad\texttt{S45-ASTRA-R-0006}
\par\noindent\textbf{Imprimé:} la convention de signe de .1
\par\noindent\textbf{Corrigé:} la convention de signe de \ref{II:eq:4-3-1}
\par\noindent The incomplete printed reference .1 refers to the alternating trace definition(4.3.1); Codex corrected restores the same reference.
\smallskip
\item[\hypertarget{apparatus:FR-0055}{55}] Page \hyperlink{printpage:fr:100}{100}\quad\texttt{S45-SOL-0386}
\par\noindent\textbf{Imprimé:} la formule (4.10.1)
\par\noindent\textbf{Corrigé:} la formule \eqref{II:eq:4-10}
\par\noindent The trace display is visibly numbered(4.10) on p.96; the printed cross-reference(4.10.1) has a spurious final component.
\smallskip
\item[\hypertarget{apparatus:FR-0056}{56}] Page \hyperlink{printpage:fr:101}{101}\quad\texttt{S45-SOL-0394}
\par\noindent\textbf{Imprimé:} cohomologie associées à un cycle
\par\noindent\textbf{Corrigé:} cohomologie associée à un cycle
\par\noindent The visible plural associées disagrees with singular classe; the corrected form uses singular agreement.
\smallskip
\item[\hypertarget{apparatus:FR-0057}{57}] Page \hyperlink{printpage:fr:102}{102}\quad\texttt{S45-SOL-0404}
\par\noindent\textbf{Imprimé:} le vérifier par $P=A[H]$
\par\noindent\textbf{Corrigé:} le vérifier pour $P=A[H]$
\par\noindent The reduction is to the case P=A[H]; the printed par is an incorrect preposition for this claim.
\smallskip
\item[\hypertarget{apparatus:FR-0058}{58}] Page \hyperlink{printpage:fr:103}{103}\quad\texttt{S45-SOL-0409}
\par\noindent\textbf{Imprimé:} l'automorphisme $\Lambda[H]$
\par\noindent\textbf{Corrigé:} l'automorphisme de $\Lambda[H]$
\par\noindent The print omits de before the algebra whose automorphism sigma is being defined.
\smallskip
\item[\hypertarget{apparatus:FR-0059}{59}] Page \hyperlink{printpage:fr:105}{105}\quad\texttt{S45-SOL-0419}
\par\noindent\textbf{Imprimé:} Le complexe de $\Lambda[H]$-modules
\par\noindent\textbf{Corrigé:} Le complexe de $A[H]$-modules
\par\noindent RΓ\_c(U,f\_*A) has coefficient ring A[H], as defined on p.104; Lambda appears only after tensoring with M.
\smallskip
\item[\hypertarget{apparatus:FR-0060}{60}] Page \hyperlink{printpage:fr:105}{105}\quad\texttt{S45-SOL-0420}
\par\noindent\textbf{Imprimé:} $\Lambda[G^+]$-modules, pour
\par\noindent\textbf{Corrigé:} $A[G^+]$-modules, pour
\par\noindent The same A[H]-complex with Frobenius action is an A[G+]-complex; the printed Lambda is inconsistent with its coefficients.
\smallskip
\item[\hypertarget{apparatus:FR-0061}{61}] Page \hyperlink{printpage:fr:105}{105}\quad\texttt{S45-SOL-0421}
\par\noindent\textbf{Imprimé:} Les formules \ref{II:5-10}
\par\noindent\textbf{Corrigé:} Les formules \ref{II:5-11}
\par\noindent The coinvariant trace formula used is proposition5.11; paragraph5.10 introduces the notation but contains no trace identity.
\smallskip
\item[\hypertarget{apparatus:FR-0062}{62}] Page \hyperlink{printpage:fr:108}{108}\quad\texttt{S45-SOL-0437}
\par\noindent\textbf{Imprimé:} rationnels de $Y_0$ vérifiant
\par\noindent\textbf{Corrigé:} rationnels de $Y_0$ vérifient
\par\noindent The finite conditional clause si les fibres ... requires the verb vérifient; the printed participle leaves it without a predicate.
\smallskip
\item[\hypertarget{apparatus:FR-0063}{63}] Page \hyperlink{printpage:fr:108}{108}\quad\texttt{S45-SOL-0438}
\par\noindent\textbf{Imprimé:} aux deux particuliers
\par\noindent\textbf{Corrigé:} aux deux cas particuliers
\par\noindent The text refers to the two special cases proved in5.1 and5.2; the printed phrase omits cas.
\smallskip
\item[\hypertarget{apparatus:FR-0064}{64}] Page \hyperlink{printpage:fr:110}{110}\quad\texttt{S45-ASTRA-R-0023}
\par\noindent\textbf{Imprimé:} Artin-Scheier
\par\noindent\textbf{Corrigé:} Artin--Schreier
\par\noindent Printed eponym misspelled.
\smallskip
\item[\hypertarget{apparatus:FR-0065}{65}] Page \hyperlink{printpage:fr:112}{112}\quad\texttt{S45-SOL-0448}
\par\noindent\textbf{Imprimé:} $A$-module projectifs
\par\noindent\textbf{Corrigé:} $A$-modules projectifs
\par\noindent Printed singular module conflicts with projectifs; plural agreement correction only.
\smallskip
\item[\hypertarget{apparatus:FR-0066}{66}] Page \hyperlink{printpage:fr:112}{112}\quad\texttt{S45-SOL-0449}
\par\noindent\textbf{Imprimé:} \(\det(1-u t,\gr_F^p(F))\)
\par\noindent\textbf{Corrigé:} \(\det(1-u t,\gr_F^p(K))\)
\par\noindent The print has F as the argument; the filtered complex is K, as stated immediately above.
\smallskip
\item[\hypertarget{apparatus:FR-0067}{67}] Page \hyperlink{printpage:fr:113}{113}\quad\texttt{S45-SOL-0450}
\par\noindent\textbf{Imprimé:} \(s(\Gamma^i K'\otimes \Gamma^i K'')\)
\par\noindent\textbf{Corrigé:} \(s(\Gamma^i K'\otimes \Gamma^j K'')\)
\par\noindent The second tensor exponent must be j in the i+j=n summand.
\smallskip
\item[\hypertarget{apparatus:FR-0068}{68}] Page \hyperlink{printpage:fr:113}{113}\quad\texttt{S45-SOL-0451}
\par\noindent\textbf{Imprimé:} déduit des complexes doubles
\par\noindent\textbf{Corrigé:} déduits des complexes doubles
\par\noindent Agreement with plural complexes simples; the print has déduit.
\smallskip
\item[\hypertarget{apparatus:FR-0069}{69}] Page \hyperlink{printpage:fr:113}{113}\quad\texttt{S45-SOL-0452}
\par\noindent\textbf{Imprimé:} $\Gamma^i K'\otimes \Gamma^i K''$, et
\par\noindent\textbf{Corrigé:} $\Gamma^i K'\otimes \Gamma^j K''$, et
\par\noindent The double-complex second tensor exponent must likewise be j.
\smallskip
\item[\hypertarget{apparatus:FR-0070}{70}] Page \hyperlink{printpage:fr:114}{114}\quad\texttt{S45-SOL-0455}
\par\noindent\textbf{Imprimé:} appliqué ou complexe $K$
\par\noindent\textbf{Corrigé:} appliqué au complexe $K$
\par\noindent The print has ou; the syntactically required preposition is au.
\smallskip
\item[\hypertarget{apparatus:FR-0071}{71}] Page \hyperlink{printpage:fr:114}{114}\quad\texttt{S45-SOL-0458}
\par\noindent\textbf{Imprimé:} \(=_{\text{(\ref{III:1-10})}}\)
\par\noindent\textbf{Corrigé:} \(=_{\text{(\ref{III:1-11})}}\)
\par\noindent The first equality uses the filtration lemma 1.10; the second uses the shifted one-term computation of 1.11.
\smallskip
\item[\hypertarget{apparatus:FR-0072}{72}] Page \hyperlink{printpage:fr:117}{117}\quad\texttt{S45-SOL-0471}
\par\noindent\textbf{Imprimé:} à degré positifs
\par\noindent\textbf{Corrigé:} à degrés positifs
\par\noindent The print omits plural s in degré.
\smallskip
\item[\hypertarget{apparatus:FR-0073}{73}] Page \hyperlink{printpage:fr:118}{118}\quad\texttt{S45-SOL-0478}
\par\noindent\textbf{Imprimé:} indenée
\par\noindent\textbf{Corrigé:} indexée
\par\noindent The enlarged scan unequivocally reads indenée; French correction to indexée. English already says indexed.
\smallskip
\item[\hypertarget{apparatus:FR-0074}{74}] Page \hyperlink{printpage:fr:118}{118}\quad\texttt{S45-SOL-0480}
\par\noindent\textbf{Imprimé:} \(\sum_n (\tr\left(F^*,\Gamma^n \Gamma(X,\sK)\right) T^n\)
\par\noindent\textbf{Corrigé:} \(\sum_n \tr\left(F^*,\Gamma^n \Gamma(X,\sK)\right) T^n\)
\par\noindent An unmatched opening parenthesis occurs immediately after the printed summation sign.
\smallskip
\item[\hypertarget{apparatus:FR-0075}{75}] Page \hyperlink{printpage:fr:118}{118}\quad\texttt{S45-SOL-0481}
\par\noindent\textbf{Imprimé:} \(\tr\left(F^*,\Gamma(\operatorname{Sym}^n(X),\Gamma_\text{ext}^n(\sK)\right)\)
\par\noindent\textbf{Corrigé:} \(\tr\left(F^*,\Gamma(\operatorname{Sym}^n(X),\Gamma_\text{ext}^n(\sK))\right)\)
\par\noindent The printed formula is one closing parenthesis short; the section functor Gamma needs closing before the trace.
\smallskip
\item[\hypertarget{apparatus:FR-0076}{76}] Page \hyperlink{printpage:fr:119}{119}\quad\texttt{S45-SOL-0483}
\par\noindent\textbf{Imprimé:} \(\tr\left(F^*,\eR\Gamma_c(\operatorname{Sym}^n(X,\Gamma_\text{ext}^n \sK)\right)\)
\par\noindent\textbf{Corrigé:} \(\tr\left(F^*,\eR\Gamma_c(\operatorname{Sym}^n(X),\Gamma_\text{ext}^n \sK)\right)\)
\par\noindent The print omits the parenthesis closing Sym\textasciicircum{}n(X).
\smallskip
\item[\hypertarget{apparatus:FR-0077}{77}] Page \hyperlink{printpage:fr:120}{120}\quad\texttt{S45-ASTRA-R-0012}
\par\noindent\textbf{Imprimé:} Théorie d'Artin-Scheier
\par\noindent\textbf{Corrigé:} Théorie d'Artin--Schreier
\par\noindent Printed eponym misspelled in section heading.
\smallskip
\item[\hypertarget{apparatus:FR-0078}{78}] Page \hyperlink{printpage:fr:120}{120}\quad\texttt{S45-ASTRA-R-0013}
\par\noindent\textbf{Imprimé:} La théorie d'Artin-Scheier
\par\noindent\textbf{Corrigé:} La théorie d'Artin--Schreier
\par\noindent Printed eponym misspelled in text.
\smallskip
\item[\hypertarget{apparatus:FR-0079}{79}] Page \hyperlink{printpage:fr:120}{120}\quad\texttt{S45-SOL-0490}
\par\noindent\textbf{Imprimé:} \(\Phi(I^n \sG')\subset I^{p^n}\Phi(\sG')\)
\par\noindent\textbf{Corrigé:} \(\Phi(I^n \sG')\subset I^{p n}\Phi(\sG')\)
\par\noindent p-semilinearity multiplies the ideal exponent by p; the printed raised n is a typo.
\smallskip
\item[\hypertarget{apparatus:FR-0080}{80}] Page \hyperlink{printpage:fr:121}{121}\quad\texttt{S45-SOL-0491}
\par\noindent\textbf{Imprimé:} une section de $x$ de
\par\noindent\textbf{Corrigé:} une section $x$ de
\par\noindent The extra de before the section name x is printed.
\smallskip
\item[\hypertarget{apparatus:FR-0081}{81}] Page \hyperlink{printpage:fr:123}{123}\quad\texttt{S45-ASTRA-R-0015}
\par\noindent\textbf{Imprimé:} La théorie d'Artin-Scheier fournit
\par\noindent\textbf{Corrigé:} La théorie d'Artin--Schreier fournit
\par\noindent Printed eponym misspelled.
\smallskip
\item[\hypertarget{apparatus:FR-0082}{82}] Page \hyperlink{printpage:fr:123}{123}\quad\texttt{S45-SOL-0501}
\par\noindent\textbf{Imprimé:} (cf. \ref{III:3-6}b)
\par\noindent\textbf{Corrigé:} (cf. \ref{III:3-3}c)
\par\noindent Section3.6 has no part b. The surjectivity argument is explicitly proved in3.3c.
\smallskip
\item[\hypertarget{apparatus:FR-0083}{83}] Page \hyperlink{printpage:fr:124}{124}\quad\texttt{S45-SOL-0506}
\par\noindent\textbf{Imprimé:} \(\ker(\Phi-1)\otimes_{\dF_q}\dF \iso\)
\par\noindent\textbf{Corrigé:} \(\ker(\Phi-1)\otimes_{\dF_p}\dF \iso\)
\par\noindent The Phi-fixed vectors form an F\_p-space. The next printed trace chain correctly uses F\_p.
\smallskip
\item[\hypertarget{apparatus:FR-0084}{84}] Page \hyperlink{printpage:fr:125}{125}\quad\texttt{S45-SOL-0509}
\par\noindent\textbf{Imprimé:} chaque membre de (3.1.1)
\par\noindent\textbf{Corrigé:} chaque membre de \eqref{III:eq:4-1-1}
\par\noindent The cited trace formula is4.1.1, not3.1.1.
\smallskip
\item[\hypertarget{apparatus:FR-0085}{85}] Page \hyperlink{printpage:fr:125}{125}\quad\texttt{S45-SOL-0512}
\par\noindent\textbf{Imprimé:} les deux membres de (3.1.1)
\par\noindent\textbf{Corrigé:} les deux membres de \eqref{III:eq:4-1-1}
\par\noindent The two trace expressions belong to4.1.1.
\smallskip
\item[\hypertarget{apparatus:FR-0086}{86}] Page \hyperlink{printpage:fr:125}{125}\quad\texttt{S45-SOL-0513}
\par\noindent\textbf{Imprimé:} Si (3.1.1) vaut
\par\noindent\textbf{Corrigé:} Si \eqref{III:eq:4-1-1} vaut
\par\noindent The reduction assumes trace formula4.1.1.
\smallskip
\item[\hypertarget{apparatus:FR-0087}{87}] Page \hyperlink{printpage:fr:126}{126}\quad\texttt{S45-ASTRA-R-0018}
\par\noindent\textbf{Imprimé:} Woodshole fixed point formula
\par\noindent\textbf{Corrigé:} Woods Hole fixed point formula
\par\noindent Correct the printed spelling of the Woods Hole name.
\smallskip
\item[\hypertarget{apparatus:FR-0088}{88}] Page \hyperlink{printpage:fr:126}{126}\quad\texttt{S45-SOL-0519}
\par\noindent\textbf{Imprimé:} appliqué à $F$ et $u$
\par\noindent\textbf{Corrigé:} appliquée à $F$ et $u$
\par\noindent Agreement with la formule requires appliquée.
\smallskip
\item[\hypertarget{apparatus:FR-0089}{89}] Page \hyperlink{printpage:fr:127}{127}\quad\texttt{S45-ASTRA-R-0019}
\par\noindent\textbf{Imprimé:} \(u^* F^*\)
\par\noindent\textbf{Corrigé:} \(u F^*\)
\par\noindent The trace is of the composition defined immediately before4.4.3; no extra star on u.
\smallskip
\item[\hypertarget{apparatus:FR-0090}{90}] Page \hyperlink{printpage:fr:127}{127}\quad\texttt{S45-SOL-0520}
\par\noindent\textbf{Imprimé:} \(\frac{\tr(u_x,\sG_x)}\)
\par\noindent\textbf{Corrigé:} \(\frac{\tr(u_x,\bar\sG_x)}\)
\par\noindent The trace must use the extended sheaf on all of barS.
\smallskip
\item[\hypertarget{apparatus:FR-0091}{91}] Page \hyperlink{printpage:fr:127}{127}\quad\texttt{S45-ASTRA-R-0020}
\par\noindent\textbf{Imprimé:} l'endomorphisme $u_x$ de $\sG_x$
\par\noindent\textbf{Corrigé:} l'endomorphisme $u_x$ de $\bar\sG_x$
\par\noindent The extended sheaf is required on the boundary; the scan omits its bar at this occurrence.
\smallskip
\item[\hypertarget{apparatus:FR-0092}{92}] Page \hyperlink{printpage:fr:133}{133}\quad\texttt{S45-SOL-0539}
\par\noindent\textbf{Imprimé:} qui nous reservirons
\par\noindent\textbf{Corrigé:} qui nous resserviront
\par\noindent The plural subject remarques requires resserviront; the printed form also omits the doubled s.
\smallskip
\item[\hypertarget{apparatus:FR-0093}{93}] Page \hyperlink{printpage:fr:137}{137}\quad\texttt{S45-ASTRA-R-0040}
\par\noindent\textbf{Imprimé:} Kunneth
\par\noindent\textbf{Corrigé:} K\"unneth
\par\noindent The print omits the diaeresis from the proper name.
\smallskip
\item[\hypertarget{apparatus:FR-0094}{94}] Page \hyperlink{printpage:fr:137}{137}\quad\texttt{S45-SOL-0552}
\par\noindent\textbf{Imprimé:} \(V^\vee\otimes V\otimes W^\vee\otimes W\to k\)
\par\noindent\textbf{Corrigé:} \(V^\vee\otimes V\otimes W^\vee\otimes W\to A\)
\par\noindent The ground field in the surrounding definition is A, including the evaluation maps1.3.1 and1.3.4; k has not been introduced.
\smallskip
\item[\hypertarget{apparatus:FR-0095}{95}] Page \hyperlink{printpage:fr:139}{139}\quad\texttt{S45-SOL-0561}
\par\noindent\textbf{Imprimé:} $D$ ne disconnecte par $X$
\par\noindent\textbf{Corrigé:} $D$ ne disconnecte pas $X$
\par\noindent The negative statement requires pas; retain the printed verb disconnecte rather than introducing Codex’s unneeded lexical replacement.
\smallskip
\item[\hypertarget{apparatus:FR-0096}{96}] Page \hyperlink{printpage:fr:141}{141}\quad\texttt{S45-SOL-0572}
\par\noindent\textbf{Imprimé:} \(\eR k^! \dZ/n=\eR k_1^! \eR i^!\)
\par\noindent\textbf{Corrigé:} \(\eR k^! \dZ/n=\eR k_1^! \eR i^!\dZ/n\)
\par\noindent The right side lacks the coefficient object, while the left is the functor applied to Z/n.
\smallskip
\item[\hypertarget{apparatus:FR-0097}{97}] Page \hyperlink{printpage:fr:142}{142}\quad\texttt{S45-SOL-0577}
\par\noindent\textbf{Imprimé:} \(\eR^{2 c+q} i^!(g^*\sF)=0\Leftrightarrow \eR^{-2 d+q}f^!(\sF)\)
\par\noindent\textbf{Corrigé:} \(\eR^{2 c+q} i^!(g^*\sF)=0\Leftrightarrow \eR^{-2 d+q}f^!(\sF)=0\)
\par\noindent The right side of the vanishing equivalence requires =0, omitted in print.
\smallskip
\item[\hypertarget{apparatus:FR-0098}{98}] Page \hyperlink{printpage:fr:143}{143}\quad\texttt{S45-SOL-0579}
\par\noindent\textbf{Imprimé:} d'où $\sI$ est
\par\noindent\textbf{Corrigé:} où $\sI$ est
\par\noindent The clause defines the ideal sheaf I, so où replaces the printed d’où.
\smallskip
\item[\hypertarget{apparatus:FR-0099}{99}] Page \hyperlink{printpage:fr:143}{143}\quad\texttt{S45-SOL-0580}
\par\noindent\textbf{Imprimé:} en terme de plongements
\par\noindent\textbf{Corrigé:} en termes de plongements
\par\noindent Printed singular terme in the expression en termes de.
\smallskip
\item[\hypertarget{apparatus:FR-0100}{100}] Page \hyperlink{printpage:fr:144}{144}\quad\texttt{S45-SOL-0585}
\par\noindent\textbf{Imprimé:} \(\h^0\left(Y,\eR^{-2 d}f^!\dZ/n\right)\)
\par\noindent\textbf{Corrigé:} \(\h^0\left(Y,\eR^{-2 d}f^!\dZ/n(-d)\right)\)
\par\noindent The preceding Hom term has twist -d, which survives into the final degree-zero cohomology group.
\smallskip
\item[\hypertarget{apparatus:FR-0101}{101}] Page \hyperlink{printpage:fr:144}{144}\quad\texttt{S45-SOL-0588}
\par\noindent\textbf{Imprimé:} \(\eR g^! \dZ/n=\dZ/n(-N)[-2 N]\)
\par\noindent\textbf{Corrigé:} \(\eR g^! \dZ/n=\dZ/n(N)[2 N]\)
\par\noindent Smooth purity for relative dimension N gives positive twist and shift, as also used on p142.
\smallskip
\item[\hypertarget{apparatus:FR-0102}{102}] Page \hyperlink{printpage:fr:145}{145}\quad\texttt{S45-ASTRA-R-0031}
\par\noindent\textbf{Imprimé:} \(\otimes \eR^{2 d} g_!\dZ/n(d)\)
\par\noindent\textbf{Corrigé:} \(\otimes \eR^{2 d} f_!\dZ/n(d)\)
\par\noindent The input pushforward must come from Y via f to match the trace relation.
\smallskip
\item[\hypertarget{apparatus:FR-0103}{103}] Page \hyperlink{printpage:fr:145}{145}\quad\texttt{S45-SOL-0593}
\par\noindent\textbf{Imprimé:} où le cup-produit peut se calculer dans $\h^*(X,\dZ/n)$.
\par\noindent\textbf{Corrigé:} où le cup-produit peut se calculer dans $\h^*(X,\dZ/n(c))$.
\par\noindent The cycle class has twist c; cup with untwisted v retains this twist.
\smallskip
\item[\hypertarget{apparatus:FR-0104}{104}] Page \hyperlink{printpage:fr:146}{146}\quad\texttt{S45-ASTRA-R-0033}
\par\noindent\textbf{Imprimé:} \PrintCite{SGA~4 XVII \S2}{\cite[XVII \S2]{sga4}}
\par\noindent\textbf{Corrigé:} \PrintCite{SGA~4 XVIII \S2}{\cite[XVIII \S2]{sga4}}
\par\noindent The printed XVII differs from the systematic trace construction of XVIII2 referred to in the continuation.
\smallskip
\item[\hypertarget{apparatus:FR-0105}{105}] Page \hyperlink{printpage:fr:146}{146}\quad\texttt{S45-SOL-0600}
\par\noindent\textbf{Imprimé:} \(u:Z\to \dA_S^d\)
\par\noindent\textbf{Corrigé:} \(u:X\to \dA_S^d\)
\par\noindent The composite ui requires u to have source X because i:Y→X.
\smallskip
\item[\hypertarget{apparatus:FR-0106}{106}] Page \hyperlink{printpage:fr:147}{147}\quad\texttt{S45-SOL-0602}
\par\noindent\textbf{Imprimé:} avec $g'$ et $g''$ lisse
\par\noindent\textbf{Corrigé:} avec $g'$ et $g''$ lisses
\par\noindent Plural agreement with the two morphisms.
\smallskip
\item[\hypertarget{apparatus:FR-0107}{107}] Page \hyperlink{printpage:fr:147}{147}\quad\texttt{S45-SOL-0603}
\par\noindent\textbf{Imprimé:} \(\h_X^{2 c}\left(Z,\dZ/n(c)\right)\)
\par\noindent\textbf{Corrigé:} \(\h_Y^{2 c}\left(X,\dZ/n(c)\right)\)
\par\noindent The lemma concerns the class of Y supported in X; Z is not an ambient scheme here.
\smallskip
\item[\hypertarget{apparatus:FR-0108}{108}] Page \hyperlink{printpage:fr:147}{147}\quad\texttt{S45-SOL-0605}
\par\noindent\textbf{Imprimé:} \(\hom_{S'}\left(\eR^{2 d'} f_!' \dZ/n,\dZ/n\right)\)
\par\noindent\textbf{Corrigé:} \(\hom_{S'}\left(\eR^{2 d'} f_!' \dZ/n(d'),\dZ/n\right)\)
\par\noindent The trace over S requires the d source twist.
\smallskip
\item[\hypertarget{apparatus:FR-0109}{109}] Page \hyperlink{printpage:fr:147}{147}\quad\texttt{S45-SOL-0606}
\par\noindent\textbf{Imprimé:} \(\eR g_!'\dZ/n(N)[2 N]\)
\par\noindent\textbf{Corrigé:} \(\eR g_!\dZ/n(N)[2 N]\)
\par\noindent The bottom row pushes forward from X to S via g; the prime is erroneous.
\smallskip
\item[\hypertarget{apparatus:FR-0110}{110}] Page \hyperlink{printpage:fr:147}{147}\quad\texttt{S45-SOL-0605}
\par\noindent\textbf{Imprimé:} \(\hom_S\left(\eR^{2 d} f_!\dZ/n,\dZ/n\right)\)
\par\noindent\textbf{Corrigé:} \(\hom_S\left(\eR^{2 d} f_!\dZ/n(d),\dZ/n\right)\)
\par\noindent The trace over S requires the d source twist.
\smallskip
\item[\hypertarget{apparatus:FR-0111}{111}] Page \hyperlink{printpage:fr:149}{149}\quad\texttt{S45-SOL-0612}
\par\noindent\textbf{Imprimé:} \(\hom_S\left(\eR^{2 d'} f_!' \dZ/n,\dZ/n\right)\)
\par\noindent\textbf{Corrigé:} \(\hom_S\left(\eR^{2 d'} f_!' \dZ/n(d'),\dZ/n\right)\)
\par\noindent The product trace requires its source twist d.
\smallskip
\item[\hypertarget{apparatus:FR-0112}{112}] Page \hyperlink{printpage:fr:149}{149}\quad\texttt{S45-SOL-0612}
\par\noindent\textbf{Imprimé:} \(\hom_S\left(\eR^{2 d}f_!\dZ/n,\dZ/n\right)\)
\par\noindent\textbf{Corrigé:} \(\hom_S\left(\eR^{2 d}f_!\dZ/n(d),\dZ/n\right)\)
\par\noindent The product trace requires its source twist d.
\smallskip
\item[\hypertarget{apparatus:FR-0113}{113}] Page \hyperlink{printpage:fr:149}{149}\quad\texttt{S45-SOL-0613}
\par\noindent\textbf{Imprimé:} pour $X'=X'$
\par\noindent\textbf{Corrigé:} pour $Y'=X'$
\par\noindent The cycle in the first factor must be all of X-prime; the printed tautology does not supply that cycle.
\smallskip
\item[\hypertarget{apparatus:FR-0114}{114}] Page \hyperlink{printpage:fr:149}{149}\quad\texttt{S45-SOL-0615}
\par\noindent\textbf{Imprimé:} \(\cl(Y')=1\in \h^0(X,\dZ/n)\)
\par\noindent\textbf{Corrigé:} \(\cl(Y')=1\in \h^0(X',\dZ/n)\)
\par\noindent The unit class belongs to the first factor X-prime.
\smallskip
\item[\hypertarget{apparatus:FR-0115}{115}] Page \hyperlink{printpage:fr:150}{150}\quad\texttt{S45-SOL-0616}
\par\noindent\textbf{Imprimé:} nour ramène ensuite
\par\noindent\textbf{Corrigé:} nous ramène ensuite
\par\noindent The scan retains the typing error nour; the first-person plural pronoun is nous.
\smallskip
\item[\hypertarget{apparatus:FR-0116}{116}] Page \hyperlink{printpage:fr:150}{150}\quad\texttt{S45-SOL-0619}
\par\noindent\textbf{Imprimé:} D'après (iv)
\par\noindent\textbf{Corrigé:} D'après (iii)
\par\noindent Theorem2.3.8 has only three assertions; its product assertion is(iii).
\smallskip
\item[\hypertarget{apparatus:FR-0117}{117}] Page \hyperlink{printpage:fr:151}{151}\quad\texttt{S45-ASTRA-R-0041}
\par\noindent\textbf{Imprimé:} Kunneth
\par\noindent\textbf{Corrigé:} K\"unneth
\par\noindent The print omits the diaeresis from the proper name.
\smallskip
\item[\hypertarget{apparatus:FR-0118}{118}] Page \hyperlink{printpage:fr:151}{151}\quad\texttt{S45-SOL-0621}
\par\noindent\textbf{Imprimé:} $2(q-N)$ $\eta_{XY}:\h^*(Y)(N-q)\to \h^*(X)$
\par\noindent\textbf{Corrigé:} $2(q-M)$ $\eta_{XY}:\h^*(Y)(M-q)\to \h^*(X)$
\par\noindent Integration over Y of dimension M shifts degree by2M and twist byM.
\smallskip
\item[\hypertarget{apparatus:FR-0119}{119}] Page \hyperlink{printpage:fr:151}{151}\quad\texttt{S45-SOL-0623}
\par\noindent\textbf{Imprimé:} trainer avec soi
\par\noindent\textbf{Corrigé:} traîner avec soi
\par\noindent The scan has no circumflex; restore it only as an explicit spelling correction.
\smallskip
\item[\hypertarget{apparatus:FR-0120}{120}] Page \hyperlink{printpage:fr:151}{151}\quad\texttt{S45-ASTRA-R-0039}
\par\noindent\textbf{Imprimé:} les twist à la Tate
\par\noindent\textbf{Corrigé:} les twists à la Tate
\par\noindent Plural article les requires the plural noun twists.
\smallskip
\item[\hypertarget{apparatus:FR-0121}{121}] Page \hyperlink{printpage:fr:152}{152}\quad\texttt{S45-SOL-0626}
\par\noindent\textbf{Imprimé:} La preuve de \ref{IV:3-4}
\par\noindent\textbf{Corrigé:} La preuve de \ref{IV:3-3}
\par\noindent The proof just given is that of Proposition3.3; Remark3.4 has no such proof.
\smallskip
\item[\hypertarget{apparatus:FR-0122}{122}] Page \hyperlink{printpage:fr:154}{154}\quad\texttt{S45-ASTRA-R-0056}
\par\noindent\textbf{Imprimé:} 1
\par\noindent\textbf{Corrigé:} 2
\par\noindent The printed local page number1 points to the introductory page; section1 actually begins on localpage2, printed155.
\smallskip
\item[\hypertarget{apparatus:FR-0123}{123}] Page \hyperlink{printpage:fr:154}{154}\quad\texttt{S45-ASTRA-R-0057}
\par\noindent\textbf{Imprimé:} 12
\par\noindent\textbf{Corrigé:} 13
\par\noindent The printed local page number12 points to the end of section3; section4 actually begins on localpage13, printed166.
\smallskip
\item[\hypertarget{apparatus:FR-0124}{124}] Page \hyperlink{printpage:fr:155}{155}\quad\texttt{S45-SOL-0634}
\par\noindent\textbf{Imprimé:} \(\hom(\sL,\Hom(\sK,\dZ/n))\)
\par\noindent\textbf{Corrigé:} \(\hom(\sL,\rHom(\sK,\dZ/n))\)
\par\noindent The derived tensor adjunction requires derived sheaf Hom. The scan omits R from this inner Hom, although the definition of D immediately above includes it.
\smallskip
\item[\hypertarget{apparatus:FR-0125}{125}] Page \hyperlink{printpage:fr:156}{156}\quad\texttt{S45-ASTRA-R-0045}
\par\noindent\textbf{Imprimé:} \(\hom(j_\ast\sF,\dZ/n)=j_\ast\sF^\vee\)
\par\noindent\textbf{Corrigé:} \(\Hom(j_\ast\sF,\dZ/n)=j_\ast\sF^\vee\)
\par\noindent The printed non-underlined Hom is an ordinary Hom group, but the right side is a sheaf; the sheaf Hom is required. The theorem states the correct sheaf identity.
\smallskip
\item[\hypertarget{apparatus:FR-0126}{126}] Page \hyperlink{printpage:fr:156}{156}\quad\texttt{S45-SOL-0639}
\par\noindent\textbf{Imprimé:} \(\eR j_\ast \sF_{\bar\eta}^I\)
\par\noindent\textbf{Corrigé:} \(\eR j_\ast\left(\sF_{\bar\eta}^I\right)^\vee\)
\par\noindent The bottom row is the dual triangle of the constant invariant sheaf; its rightmost term must carry the dual, as the next-page exact sequence confirms.
\smallskip
\item[\hypertarget{apparatus:FR-0127}{127}] Page \hyperlink{printpage:fr:158}{158}\quad\texttt{S45-ASTRA-R-0048}
\par\noindent\textbf{Imprimé:} \(\rhom(f_\ast\dZ/n,\sK_Y)\)
\par\noindent\textbf{Corrigé:} \(\rHom(f_\ast\dZ/n,\sK_Y)\)
\par\noindent The printed ordinary derived Hom is a complex of groups, while the equality concerns complexes of sheaves on Y; derived sheaf Hom is required.
\smallskip
\item[\hypertarget{apparatus:FR-0128}{128}] Page \hyperlink{printpage:fr:158}{158}\quad\texttt{S45-SOL-0642}
\par\noindent\textbf{Imprimé:} \(f_\ast\dZ/n\)
\par\noindent\textbf{Corrigé:} \(f_\ast\dZ/n(1)[2]\)
\par\noindent The dualizing complex K\_Y is Z/n(1)[2], so the final sheaf must carry the same twist and shift. The scan omits both.
\smallskip
\item[\hypertarget{apparatus:FR-0129}{129}] Page \hyperlink{printpage:fr:162}{162}\quad\texttt{S45-SOL-0655}
\par\noindent\textbf{Imprimé:} \eqref{V:eq:2-2-1}
\par\noindent\textbf{Corrigé:} \eqref{V:eq:3-2-1}
\par\noindent The triangle is defined by exact sequence3.2.1 immediately above, not theorem2.2.1.
\smallskip
\item[\hypertarget{apparatus:FR-0130}{130}] Page \hyperlink{printpage:fr:162}{162}\quad\texttt{S45-SOL-0656}
\par\noindent\textbf{Imprimé:} \(\eR \pr_{2\ast}\sG\)
\par\noindent\textbf{Corrigé:} \(\eR \pr_{1\ast}\sG\)
\par\noindent Both G terms use the same direct-image functor Rpr\_1*.
\smallskip
\item[\hypertarget{apparatus:FR-0131}{131}] Page \hyperlink{printpage:fr:162}{162}\quad\texttt{S45-SOL-0656}
\par\noindent\textbf{Imprimé:} \(\eR \pr_{2\ast}\sG\)
\par\noindent\textbf{Corrigé:} \(\eR \pr_{1\ast}\sG\)
\par\noindent Both G terms use the same direct-image functor Rpr\_1*.
\smallskip
\item[\hypertarget{apparatus:FR-0132}{132}] Page \hyperlink{printpage:fr:163}{163}\quad\texttt{S45-SOL-0657}
\par\noindent\textbf{Imprimé:} \(3.2.2\)
\par\noindent\textbf{Corrigé:} \(3.2.3\)
\par\noindent The print duplicates the number3.2.2 already assigned to the distinguished triangle. The next diagram is3.2.3; preserve the duplicated printed number diplomatically.
\smallskip
\item[\hypertarget{apparatus:FR-0133}{133}] Page \hyperlink{printpage:fr:163}{163}\quad\texttt{S45-SOL-0659}
\par\noindent\textbf{Imprimé:} une suite exactes de complexes
\par\noindent\textbf{Corrigé:} une suite exacte de complexes
\par\noindent Singular suite requires singular exacte; the plural s is visible in the scan.
\smallskip
\item[\hypertarget{apparatus:FR-0134}{134}] Page \hyperlink{printpage:fr:165}{165}\quad\texttt{S45-SOL-0664}
\par\noindent\textbf{Imprimé:} \(u'=c^{1,1}\)
\par\noindent\textbf{Corrigé:} \(u'=-c^{1,1}\)
\par\noindent Proposition3.2 gives c\textasciicircum{}\{1,1\}=-u. The identification defining u-prime therefore gives u-prime=-c\textasciicircum{}\{1,1\}, consistent with the following trace equality.
\smallskip
\item[\hypertarget{apparatus:FR-0135}{135}] Page \hyperlink{printpage:fr:168}{168}\quad\texttt{S45-ASTRA-S-0002}
\par\noindent\textbf{Imprimé:} généralisée
\par\noindent\textbf{Corrigé:} généralisées
\par\noindent Agreement with plural Sommes; the scan has -ée.
\smallskip
\item[\hypertarget{apparatus:FR-0136}{136}] Page \hyperlink{printpage:fr:171}{171}\quad\texttt{S45-SOL-0682}
\par\noindent\textbf{Imprimé:} objects
\par\noindent\textbf{Corrigé:} objets
\par\noindent French spelling: objets has no c; the printed reading is directly visible.
\smallskip
\item[\hypertarget{apparatus:FR-0137}{137}] Page \hyperlink{printpage:fr:173}{173}\quad\texttt{S45-SOL-0685}
\par\noindent\textbf{Imprimé:} joint
\par\noindent\textbf{Corrigé:} jointe
\par\noindent Agreement of the participle with la définition.
\smallskip
\item[\hypertarget{apparatus:FR-0138}{138}] Page \hyperlink{printpage:fr:173}{173}\quad\texttt{S45-SOL-0690}
\par\noindent\textbf{Imprimé:} \(\sF(\chi,f)\)
\par\noindent\textbf{Corrigé:} \(\sF(\chi,f_0)\)
\par\noindent The morphism in the explicitly cited 1.7 d) is f\_0.
\smallskip
\item[\hypertarget{apparatus:FR-0139}{139}] Page \hyperlink{printpage:fr:174}{174}\quad\texttt{S45-SOL-0692}
\par\noindent\textbf{Imprimé:} Pour $k$ est une extension
\par\noindent\textbf{Corrigé:} Pour $k$ une extension
\par\noindent The finite-verb est is incompatible with the printed introductory Pour; remove the redundant verb.
\smallskip
\item[\hypertarget{apparatus:FR-0140}{140}] Page \hyperlink{printpage:fr:176}{176}\quad\texttt{S45-SOL-0696}
\par\noindent\textbf{Imprimé:} \(\h_c^p(X,\dQ_\ell)\)
\par\noindent\textbf{Corrigé:} \(\h_c^i(X,\dQ_\ell)\)
\par\noindent The preceding isomorphism compares degrees i and i-2.
\smallskip
\item[\hypertarget{apparatus:FR-0141}{141}] Page \hyperlink{printpage:fr:176}{176}\quad\texttt{S45-SOL-0697}
\par\noindent\textbf{Imprimé:} \(\sum (-1)^i \tr(F^\ast,\h^i(X,\dQ_\ell))\)
\par\noindent\textbf{Corrigé:} \(\sum (-1)^i \tr(F^\ast,\h_c^i(X,\dQ_\ell))\)
\par\noindent The point-counting trace formula uses compact support for nonproper X.
\smallskip
\item[\hypertarget{apparatus:FR-0142}{142}] Page \hyperlink{printpage:fr:178}{178}\quad\texttt{S45-SOL-0707}
\par\noindent\textbf{Imprimé:} $\h_c^i(X,\sF)$ et $\h^{2n-i}(X,\sF^\vee)(2n)$
\par\noindent\textbf{Corrigé:} $\h_c^i(S,\sF)$ et $\h^{2n-i}(S,\sF^\vee)(n)$
\par\noindent The scheme in the hypothesis is S and the Poincaré twist in dimension n is n.
\smallskip
\item[\hypertarget{apparatus:FR-0143}{143}] Page \hyperlink{printpage:fr:178}{178}\quad\texttt{S45-SOL-0710}
\par\noindent\textbf{Imprimé:} \(\h^0(U,\sF^\vee)(2n)\)
\par\noindent\textbf{Corrigé:} \(\h^0(U,\sF^\vee)(n)\)
\par\noindent Poincaré duality in dimension n supplies twist n.
\smallskip
\item[\hypertarget{apparatus:FR-0144}{144}] Page \hyperlink{printpage:fr:179}{179}\quad\texttt{S45-SOL-0711}
\par\noindent\textbf{Imprimé:} \((\sF_u)_{\pi_1(U,u)}(-2n)\)
\par\noindent\textbf{Corrigé:} \((\sF_u)_{\pi_1(U,u)}(-n)\)
\par\noindent Dualizing the dimension-n Tate twist gives -n.
\smallskip
\item[\hypertarget{apparatus:FR-0145}{145}] Page \hyperlink{printpage:fr:179}{179}\quad\texttt{S45-SOL-0719}
\par\noindent\textbf{Imprimé:} un ouvert de $X$
\par\noindent\textbf{Corrigé:} un ouvert de $\bar X$
\par\noindent The open subset is X inside the proper ambient bar X.
\smallskip
\item[\hypertarget{apparatus:FR-0146}{146}] Page \hyperlink{printpage:fr:180}{180}\quad\texttt{S45-SOL-0721}
\par\noindent\textbf{Imprimé:} (cf. \ref{VI:1-17})
\par\noindent\textbf{Corrigé:} (cf. \ref{VI:1-19})
\par\noindent The comparison criterion immediately above is Proposition 1.19.
\smallskip
\item[\hypertarget{apparatus:FR-0147}{147}] Page \hyperlink{printpage:fr:180}{180}\quad\texttt{S45-SOL-0722}
\par\noindent\textbf{Imprimé:} Appliquons \ref{VI:1-15}
\par\noindent\textbf{Corrigé:} Appliquons \ref{VI:1-17}
\par\noindent The argument invokes the weight bound of Theorem 1.17.
\smallskip
\item[\hypertarget{apparatus:FR-0148}{148}] Page \hyperlink{printpage:fr:181}{181}\quad\texttt{S45-SOL-0728}
\par\noindent\textbf{Imprimé:} $\rho(g^{-1} \gamma g)$ de $E_\gamma^n$
\par\noindent\textbf{Corrigé:} $\rho(g^{-1} \gamma g)$ de $E_\lambda^n$
\par\noindent The coefficient field specified in the representation is E\_lambda; gamma labels the rational point.
\smallskip
\item[\hypertarget{apparatus:FR-0149}{149}] Page \hyperlink{printpage:fr:183}{183}\quad\texttt{S45-SOL-0732}
\par\noindent\textbf{Imprimé:} il portent
\par\noindent\textbf{Corrigé:} ils portent
\par\noindent Plural subject and plural verb require ils.
\smallskip
\item[\hypertarget{apparatus:FR-0150}{150}] Page \hyperlink{printpage:fr:184}{184}\quad\texttt{S45-SOL-0736}
\par\noindent\textbf{Imprimé:} \(\tr(F^\ast,\h^\bullet(X,\sF))\)
\par\noindent\textbf{Corrigé:} \(\tr(F^\ast,\h_c^\bullet(X,\sF))\)
\par\noindent The trace formula requires compact supports without properness.
\smallskip
\item[\hypertarget{apparatus:FR-0151}{151}] Page \hyperlink{printpage:fr:184}{184}\quad\texttt{S45-SOL-0737}
\par\noindent\textbf{Imprimé:} \(\tr(F^\ast, \h^\bullet(Y,\eR^q f_! \sF))\)
\par\noindent\textbf{Corrigé:} \(\tr(F^\ast,\h_c^\bullet(Y,\eR^q f_! \sF))\)
\par\noindent The trace formula on Y is a compact-support trace formula.
\smallskip
\item[\hypertarget{apparatus:FR-0152}{152}] Page \hyperlink{printpage:fr:185}{185}\quad\texttt{S45-SOL-0738}
\par\noindent\textbf{Imprimé:} \(\tr(F^\ast,\h^\bullet(Y,\eR^\bullet f_! \sF))\)
\par\noindent\textbf{Corrigé:} \(\tr(F^\ast,\h_c^\bullet(Y,\eR^\bullet f_! \sF))\)
\par\noindent Alternating the compact-support Leray sequence preserves compact supports.
\smallskip
\item[\hypertarget{apparatus:FR-0153}{153}] Page \hyperlink{printpage:fr:186}{186}\quad\texttt{S45-SOL-0745}
\par\noindent\textbf{Imprimé:} n'était pas exclus
\par\noindent\textbf{Corrigé:} n'était pas exclu
\par\noindent Agreement with the singular case J=emptyset.
\smallskip
\item[\hypertarget{apparatus:FR-0154}{154}] Page \hyperlink{printpage:fr:187}{187}\quad\texttt{S45-SOL-0748}
\par\noindent\textbf{Imprimé:} \((\chi^{-1}(g)\)
\par\noindent\textbf{Corrigé:} \(\chi^{-1}(g)\)
\par\noindent Delete the unmatched opening parenthesis.
\smallskip
\item[\hypertarget{apparatus:FR-0155}{155}] Page \hyperlink{printpage:fr:187}{187}\quad\texttt{S45-SOL-0749}
\par\noindent\textbf{Imprimé:} \(\h^\bullet(G,\sF(\chi))\)
\par\noindent\textbf{Corrigé:} \(\h_c^\bullet(G,\sF(\chi))\)
\par\noindent rho\_H was defined on compactly supported cohomology.
\smallskip
\item[\hypertarget{apparatus:FR-0156}{156}] Page \hyperlink{printpage:fr:188}{188}\quad\texttt{S45-SOL-0750}
\par\noindent\textbf{Imprimé:} \(\h^p(X,\sF)\)
\par\noindent\textbf{Corrigé:} \(\h_c^p(X,\sF)\)
\par\noindent The fibre of R\textasciicircum{}p pr\_\{1!\} has compact supports.
\smallskip
\item[\hypertarget{apparatus:FR-0157}{157}] Page \hyperlink{printpage:fr:189}{189}\quad\texttt{S45-SOL-0757}
\par\noindent\textbf{Imprimé:} \(\h^{2-i}(\bar X,j_\ast(\sF^\vee(1))\)
\par\noindent\textbf{Corrigé:} \(\h^{2-i}(\bar X,j_\ast(\sF^\vee(1)))\)
\par\noindent Balance the parentheses of the cohomology group.
\smallskip
\item[\hypertarget{apparatus:FR-0158}{158}] Page \hyperlink{printpage:fr:190}{190}\quad\texttt{S45-SOL-0760}
\par\noindent\textbf{Imprimé:} \(\psi(a^p-a)=0\)
\par\noindent\textbf{Corrigé:} \(\psi(a^p-a)=1\)
\par\noindent The defining exponential has value 1 when the trace is zero.
\smallskip
\item[\hypertarget{apparatus:FR-0159}{159}] Page \hyperlink{printpage:fr:190}{190}\quad\texttt{S45-SOL-0763}
\par\noindent\textbf{Imprimé:} Cette somme est nulle pour $f$ de la forme $g^p-g$.
\par\noindent\textbf{Corrigé:} Pour $f$ de la forme $g^p-g$, cette somme est égale à $\#\{x\in X_0(\dF_q)\mid f(x)\ne\infty\}$.
\par\noindent Every summand in this case is 1.
\smallskip
\item[\hypertarget{apparatus:FR-0160}{160}] Page \hyperlink{printpage:fr:190}{190}\quad\texttt{S45-SOL-0765}
\par\noindent\textbf{Imprimé:} pour tout point fermé $s$
\par\noindent\textbf{Corrigé:} pour tout point fermé $x$
\par\noindent Subsequent notation uniformly uses x.
\smallskip
\item[\hypertarget{apparatus:FR-0161}{161}] Page \hyperlink{printpage:fr:190}{190}\quad\texttt{S45-SOL-0766}
\par\noindent\textbf{Imprimé:} pole de $f$
\par\noindent\textbf{Corrigé:} pôle de $f$
\par\noindent French pôle requires its circumflex; the print omits it.
\smallskip
\item[\hypertarget{apparatus:FR-0162}{162}] Page \hyperlink{printpage:fr:191}{191}\quad\texttt{S45-SOL-0771}
\par\noindent\textbf{Imprimé:} celle du $\dF_q$-torseur
\par\noindent\textbf{Corrigé:} celle du $\dF_p$-torseur
\par\noindent The kernel of the Artin–Schreier polynomial T\textasciicircum{}p-T is F\_p.
\smallskip
\item[\hypertarget{apparatus:FR-0163}{163}] Page \hyperlink{printpage:fr:192}{192}\quad\texttt{S45-SOL-0773}
\par\noindent\textbf{Imprimé:} $\dF_q$-torseur sur $\spec(\dF_q)$
\par\noindent\textbf{Corrigé:} $\dF_p$-torseur sur $\spec(\dF_q)$
\par\noindent The structure group is F\_p whereas the base remains Spec(F\_q).
\smallskip
\item[\hypertarget{apparatus:FR-0164}{164}] Page \hyperlink{printpage:fr:192}{192}\quad\texttt{S45-SOL-0780}
\par\noindent\textbf{Imprimé:} dualité parfaite à valeur dans $E_\lambda(1)$
\par\noindent\textbf{Corrigé:} dualité parfaite à valeur dans $E_\lambda(-1)$
\par\noindent Poincaré duality pairs degree-one groups into H\textasciicircum{}2(E\_lambda)=E\_lambda(-1).
\smallskip
\item[\hypertarget{apparatus:FR-0165}{165}] Page \hyperlink{printpage:fr:192}{192}\quad\texttt{S45-SOL-0781}
\par\noindent\textbf{Imprimé:} $\sigma^\ast$ agit trivialement sur $E_\lambda(1)$
\par\noindent\textbf{Corrigé:} $\sigma^\ast$ agit trivialement sur $E_\lambda(-1)$
\par\noindent Maintain the corrected Poincaré-duality target.
\smallskip
\item[\hypertarget{apparatus:FR-0166}{166}] Page \hyperlink{printpage:fr:193}{193}\quad\texttt{S45-SOL-0784}
\par\noindent\textbf{Imprimé:} à $n$ variable
\par\noindent\textbf{Corrigé:} à $n$ variables
\par\noindent Plural agreement with n variables.
\smallskip
\item[\hypertarget{apparatus:FR-0167}{167}] Page \hyperlink{printpage:fr:193}{193}\quad\texttt{S45-SOL-0786}
\par\noindent\textbf{Imprimé:} $r\geqslant i$
\par\noindent\textbf{Corrigé:} $r\geqslant 1$
\par\noindent The lower bound in the induction is the integer 1.
\smallskip
\item[\hypertarget{apparatus:FR-0168}{168}] Page \hyperlink{printpage:fr:193}{193}\quad\texttt{S45-SOL-0787}
\par\noindent\textbf{Imprimé:} par tout les points
\par\noindent\textbf{Corrigé:} par tous les points
\par\noindent Plural agreement of tous with points.
\smallskip
\item[\hypertarget{apparatus:FR-0169}{169}] Page \hyperlink{printpage:fr:193}{193}\quad\texttt{S45-SOL-0789}
\par\noindent\textbf{Imprimé:} hyperplans rationnel
\par\noindent\textbf{Corrigé:} hyperplans rationnels
\par\noindent Plural agreement with hyperplans.
\smallskip
\item[\hypertarget{apparatus:FR-0170}{170}] Page \hyperlink{printpage:fr:194}{194}\quad\texttt{S45-SOL-0792}
\par\noindent\textbf{Imprimé:} {\footnotesize\ttfamily\raggedright \textbackslash{}\allowbreak{}sum\_\allowbreak{}\{\allowbreak{}x\_\allowbreak{}2,\allowbreak{}\textbackslash{}\allowbreak{}dots,\allowbreak{}x\_\allowbreak{}n\}\allowbreak{} \textbackslash{}\allowbreak{}sum\_\allowbreak{}\{\allowbreak{}x\_\allowbreak{}1\}\allowbreak{} \textbackslash{}\allowbreak{}psi P(\allowbreak{}x\_\allowbreak{}1,\allowbreak{}\textbackslash{}\allowbreak{}dots,\allowbreak{}x\_\allowbreak{}n\textbackslash{}\allowbreak{}right|}
\par\noindent\textbf{Corrigé:} {\footnotesize\ttfamily\raggedright \textbackslash{}\allowbreak{}sum\_\allowbreak{}\{\allowbreak{}x\_\allowbreak{}2,\allowbreak{}\textbackslash{}\allowbreak{}dots,\allowbreak{}x\_\allowbreak{}n\}\allowbreak{} \textbackslash{}\allowbreak{}sum\_\allowbreak{}\{\allowbreak{}x\_\allowbreak{}1\}\allowbreak{} \textbackslash{}\allowbreak{}psi P(\allowbreak{}x\_\allowbreak{}1,\allowbreak{}\textbackslash{}\allowbreak{}dots,\allowbreak{}x\_\allowbreak{}n)\allowbreak{}\textbackslash{}\allowbreak{}right|}
\par\noindent Close the argument of P before the absolute-value delimiter.
\smallskip
\item[\hypertarget{apparatus:FR-0171}{171}] Page \hyperlink{printpage:fr:195}{195}\quad\texttt{S45-SOL-0794}
\par\noindent\textbf{Imprimé:} {\footnotesize\ttfamily\raggedright \textbackslash{}\allowbreak{}left.}
\par\noindent\textbf{Corrigé:} {\footnotesize\ttfamily\raggedright \textbackslash{}\allowbreak{}left|}
\par\noindent Supply the missing opening absolute-value delimiter.
\smallskip
\item[\hypertarget{apparatus:FR-0172}{172}] Page \hyperlink{printpage:fr:195}{195}\quad\texttt{S45-SOL-0795}
\par\noindent\textbf{Imprimé:} \(\sum_{S=\lambda}\)
\par\noindent\textbf{Corrigé:} \(\sum_{S_1=\lambda}\)
\par\noindent The linear form in the argument is S\_1.
\smallskip
\item[\hypertarget{apparatus:FR-0173}{173}] Page \hyperlink{printpage:fr:195}{195}\quad\texttt{S45-SOL-0796}
\par\noindent\textbf{Imprimé:} \(\sum_{S(X_2,\dots,X_n)\ne\lambda}\)
\par\noindent\textbf{Corrigé:} \(\sum_{S_1(X_2,\dots,X_n)\ne\lambda}\)
\par\noindent The same linear form is S\_1.
\smallskip
\item[\hypertarget{apparatus:FR-0174}{174}] Page \hyperlink{printpage:fr:195}{195}\quad\texttt{S45-SOL-0797}
\par\noindent\textbf{Imprimé:} \(< (d-2) q^{n-1/2} + q^{n-1}\)
\par\noindent\textbf{Corrigé:} \(\leqslant (d-2) q^{n-1/2} + q^{n-1}\)
\par\noindent Equality occurs at this intermediate step when d=2.
\smallskip
\item[\hypertarget{apparatus:FR-0175}{175}] Page \hyperlink{printpage:fr:196}{196}\quad\texttt{S45-SOL-0799}
\par\noindent\textbf{Imprimé:} \(\tr(F^\ast,\h_c^\ast(\dG_m,\sF(\psi \chi^{-1}))\)
\par\noindent\textbf{Corrigé:} \(\tr(F^\ast,\h_c^\ast(\dG_m,\sF(\psi \chi^{-1})))\)
\par\noindent Close the trace expression.
\smallskip
\item[\hypertarget{apparatus:FR-0176}{176}] Page \hyperlink{printpage:fr:197}{197}\quad\texttt{S45-SOL-0805}
\par\noindent\textbf{Imprimé:} \(|\tau(\chi,\psi)|=q\)
\par\noindent\textbf{Corrigé:} \(|\tau(\chi,\psi)|=q^{1/2}\)
\par\noindent The product with the conjugate is the square of the absolute value.
\smallskip
\item[\hypertarget{apparatus:FR-0177}{177}] Page \hyperlink{printpage:fr:197}{197}\quad\texttt{S45-SOL-0806}
\par\noindent\textbf{Imprimé:} un produits de corps
\par\noindent\textbf{Corrigé:} un produit de corps
\par\noindent Singular agreement.
\smallskip
\item[\hypertarget{apparatus:FR-0178}{178}] Page \hyperlink{printpage:fr:198}{198}\quad\texttt{S45-SOL-0810}
\par\noindent\textbf{Imprimé:} les schémas déduits de $V$ suivant
\par\noindent\textbf{Corrigé:} les schémas déduits de $V$ suivants
\par\noindent Plural agreement of the adjective with schémas.
\smallskip
\item[\hypertarget{apparatus:FR-0179}{179}] Page \hyperlink{printpage:fr:198}{198}\quad\texttt{S45-SOL-0811}
\par\noindent\textbf{Imprimé:} $Q$ $P$ sur
\par\noindent\textbf{Corrigé:} $Q$ et $P$ sur
\par\noindent Restore the missing conjunction joining the two projective spaces.
\smallskip
\item[\hypertarget{apparatus:FR-0180}{180}] Page \hyperlink{printpage:fr:200}{200}\quad\texttt{S45-SOL-0817}
\par\noindent\textbf{Imprimé:} indénée par $I$
\par\noindent\textbf{Corrigé:} indexée par $I$
\par\noindent Correct the malformed French word.
\smallskip
\item[\hypertarget{apparatus:FR-0181}{181}] Page \hyperlink{printpage:fr:200}{200}\quad\texttt{S45-SOL-0819}
\par\noindent\textbf{Imprimé:} \(\gal(\bar\dF/F)\)
\par\noindent\textbf{Corrigé:} \(\gal(\dF/\dF_q)\)
\par\noindent Arithmetic Galois action over the finite base field.
\smallskip
\item[\hypertarget{apparatus:FR-0182}{182}] Page \hyperlink{printpage:fr:200}{200}\quad\texttt{S45-SOL-0820}
\par\noindent\textbf{Imprimé:} un plongement $\omega$ de $k$ dans $\dF_q$
\par\noindent\textbf{Corrigé:} un plongement $\omega$ de $k$ dans $\dF$
\par\noindent A general finite extension k embeds into the algebraic closure, not the base field.
\smallskip
\item[\hypertarget{apparatus:FR-0183}{183}] Page \hyperlink{printpage:fr:201}{201}\quad\texttt{S45-SOL-0826}
\par\noindent\textbf{Imprimé:} une famille de caractère
\par\noindent\textbf{Corrigé:} une famille de caractères
\par\noindent The indexed objects are characters.
\smallskip
\item[\hypertarget{apparatus:FR-0184}{184}] Page \hyperlink{printpage:fr:202}{202}\quad\texttt{S45-SOL-0828}
\par\noindent\textbf{Imprimé:} \(\tr_{\dF_q/k}\)
\par\noindent\textbf{Corrigé:} \(\tr_{k/\dF_q}\)
\par\noindent Reverse the incorrectly ordered source and target of the trace.
\smallskip
\item[\hypertarget{apparatus:FR-0185}{185}] Page \hyperlink{printpage:fr:202}{202}\quad\texttt{S45-SOL-0830}
\par\noindent\textbf{Imprimé:} \(N^\ast\sF(\chi) = \bigotimes_i \pr_i^\ast\sF(\psi)\)
\par\noindent\textbf{Corrigé:} \(N^\ast\sF(\chi) = \bigotimes_i \pr_i^\ast\sF(\chi)\)
\par\noindent The norm pullback factors into copies of the multiplicative-character sheaf.
\smallskip
\item[\hypertarget{apparatus:FR-0186}{186}] Page \hyperlink{printpage:fr:202}{202}\quad\texttt{S45-SOL-0832}
\par\noindent\textbf{Imprimé:} du cardinal $n$ de $I$
\par\noindent\textbf{Corrigé:} du cardinal $N$ de $I$
\par\noindent The cardinality of I is N.
\smallskip
\item[\hypertarget{apparatus:FR-0187}{187}] Page \hyperlink{printpage:fr:202}{202}\quad\texttt{S45-SOL-0833}
\par\noindent\textbf{Imprimé:} {\footnotesize\ttfamily\raggedright \textbackslash{}\allowbreak{}h\_\allowbreak{}c\textasciicircum{}\allowbreak{}N\textbackslash{}\allowbreak{}left(\allowbreak{}V,\allowbreak{}}
\par\noindent\textbf{Corrigé:} {\footnotesize\ttfamily\raggedright \textbackslash{}\allowbreak{}h\_\allowbreak{}c\textasciicircum{}\allowbreak{}N\textbackslash{}\allowbreak{}left(\allowbreak{}V\textasciicircum{}\allowbreak{}\textbackslash{}\allowbreak{}ast,\allowbreak{}}
\par\noindent The character sheaf is on the unit torus.
\smallskip
\item[\hypertarget{apparatus:FR-0188}{188}] Page \hyperlink{printpage:fr:202}{202}\quad\texttt{S45-SOL-0834}
\par\noindent\textbf{Imprimé:} \(\h_c^1(\dG_m,\sF(\psi\chi^{-1})) \otimes_\dZ\)
\par\noindent\textbf{Corrigé:} \(\h_c^1(\dG_m,\sF(\psi\chi^{-1}))^{\otimes N} \otimes_\dZ\)
\par\noindent Künneth requires N copies of the degree-one cohomology.
\smallskip
\item[\hypertarget{apparatus:FR-0189}{189}] Page \hyperlink{printpage:fr:203}{203}\quad\texttt{S45-SOL-0838}
\par\noindent\textbf{Imprimé:} copes
\par\noindent\textbf{Corrigé:} copies
\par\noindent The restored paragraph visibly says copes. It is a misspelling of copies.
\smallskip
\item[\hypertarget{apparatus:FR-0190}{190}] Page \hyperlink{printpage:fr:203}{203}\quad\texttt{S45-SOL-0839}
\par\noindent\textbf{Imprimé:} \((k^\ast)^{N+1}\)
\par\noindent\textbf{Corrigé:} \((\dF_q^\ast)^{N+1}\)
\par\noindent The coordinate tuple is printed as an element of k\textasciicircum{}\{* N+1\}, but each coordinate belongs to F\_q\textasciicircum{}* because k is the product of N+1 copies of F\_q. Correct the coordinate field.
\smallskip
\item[\hypertarget{apparatus:FR-0191}{191}] Page \hyperlink{printpage:fr:203}{203}\quad\texttt{S45-SOL-0840}
\par\noindent\textbf{Imprimé:} \(n\)
\par\noindent\textbf{Corrigé:} \(N\)
\par\noindent The upper limit in the trace-zero condition is printed as lowercase n. The coordinates are indexed from 1 to N, so the bound is N.
\smallskip
\item[\hypertarget{apparatus:FR-0192}{192}] Page \hyperlink{printpage:fr:203}{203}\quad\texttt{S45-SOL-0841}
\par\noindent\textbf{Imprimé:} \(\chi_n^{-1}(x_n)\)
\par\noindent\textbf{Corrigé:} \(\chi_N^{-1}(x_N)\)
\par\noindent The final factor in (4.14.2) visibly has lowercase n in both chi\_n and x\_n, although the summation uses N coordinates. Correct both indices to N.
\smallskip
\item[\hypertarget{apparatus:FR-0193}{193}] Page \hyperlink{printpage:fr:204}{204}\quad\texttt{S45-SOL-0843}
\par\noindent\textbf{Imprimé:} \(q\sum_{x\in k^\ast/\dF_q^\ast} \chi(x)^{-1}\)
\par\noindent\textbf{Corrigé:} \(q\sum_{\substack{x\in k^\ast/\dF_q^\ast\\\tr(x)=0}} \chi(x)^{-1}\)
\par\noindent The first term comes from the trace-zero indicator.
\smallskip
\item[\hypertarget{apparatus:FR-0194}{194}] Page \hyperlink{printpage:fr:205}{205}\quad\texttt{S45-SOL-0847}
\par\noindent\textbf{Imprimé:} \(\sK_m((\chi_i)_{i\in I})\)
\par\noindent\textbf{Corrigé:} \(\sK_n((\chi_i)_{i\in I})\)
\par\noindent The Kummer exponent of the given characters is n.
\smallskip
\item[\hypertarget{apparatus:FR-0195}{195}] Page \hyperlink{printpage:fr:205}{205}\quad\texttt{S45-SOL-0847}
\par\noindent\textbf{Imprimé:} \(\sK_m((\chi_i)_{i\in I})\)
\par\noindent\textbf{Corrigé:} \(\sK_n((\chi_i)_{i\in I})\)
\par\noindent The Kummer exponent of the given characters is n.
\smallskip
\item[\hypertarget{apparatus:FR-0196}{196}] Page \hyperlink{printpage:fr:205}{205}\quad\texttt{S45-SOL-0849}
\par\noindent\textbf{Imprimé:} $T_0 V_0$ s'annule
\par\noindent\textbf{Corrigé:} $T_0 v_0$ s'annule
\par\noindent Composition with the blow-up morphism v\_0.
\smallskip
\item[\hypertarget{apparatus:FR-0197}{197}] Page \hyperlink{printpage:fr:205}{205}\quad\texttt{S45-SOL-0850}
\par\noindent\textbf{Imprimé:} sur $\pi^{-1}(Q_0)$
\par\noindent\textbf{Corrigé:} sur $\bar\pi^{-1}(Q_0)$
\par\noindent Use the projection of the full line bundle.
\smallskip
\item[\hypertarget{apparatus:FR-0198}{198}] Page \hyperlink{printpage:fr:205}{205}\quad\texttt{S45-SOL-0851}
\par\noindent\textbf{Imprimé:} $D=\pi^{-1}(x)$
\par\noindent\textbf{Corrigé:} $D=\bar\pi^{-1}(x)$
\par\noindent The full line is the fibre of bar pi.
\smallskip
\item[\hypertarget{apparatus:FR-0199}{199}] Page \hyperlink{printpage:fr:205}{205}\quad\texttt{S45-SOL-0852}
\par\noindent\textbf{Imprimé:} est envoyé
\par\noindent\textbf{Corrigé:} est envoyée
\par\noindent Agreement with la droite.
\smallskip
\item[\hypertarget{apparatus:FR-0200}{200}] Page \hyperlink{printpage:fr:206}{206}\quad\texttt{S45-SOL-0856}
\par\noindent\textbf{Imprimé:} \(N-1\)
\par\noindent\textbf{Corrigé:} \(N-2\)
\par\noindent With the revised rank N of M, the projective trace-zero locus has dimension N-2.
\smallskip
\item[\hypertarget{apparatus:FR-0201}{201}] Page \hyperlink{printpage:fr:206}{206}\quad\texttt{S45-SOL-0856}
\par\noindent\textbf{Imprimé:} \(N-1\)
\par\noindent\textbf{Corrigé:} \(N-2\)
\par\noindent With the revised rank N of M, the projective trace-zero locus has dimension N-2.
\smallskip
\item[\hypertarget{apparatus:FR-0202}{202}] Page \hyperlink{printpage:fr:206}{206}\quad\texttt{S45-SOL-0856}
\par\noindent\textbf{Imprimé:} \(N-1\)
\par\noindent\textbf{Corrigé:} \(N-2\)
\par\noindent With the revised rank N of M, the projective trace-zero locus has dimension N-2.
\smallskip
\item[\hypertarget{apparatus:FR-0203}{203}] Page \hyperlink{printpage:fr:207}{207}\quad\texttt{S45-SOL-0858}
\par\noindent\textbf{Imprimé:} modérée a l'infini
\par\noindent\textbf{Corrigé:} modérée à l'infini
\par\noindent The preposition requires the grave accent.
\smallskip
\item[\hypertarget{apparatus:FR-0204}{204}] Page \hyperlink{printpage:fr:208}{208}\quad\texttt{S45-SOL-0859}
\par\noindent\textbf{Imprimé:} et $K$ un corps
\par\noindent\textbf{Corrigé:} et $k$ un corps
\par\noindent The field is denoted k throughout the subsequent definitions.
\smallskip
\item[\hypertarget{apparatus:FR-0205}{205}] Page \hyperlink{printpage:fr:208}{208}\quad\texttt{S45-SOL-0860}
\par\noindent\textbf{Imprimé:} diverses façon équivalentes
\par\noindent\textbf{Corrigé:} diverses façons équivalentes
\par\noindent Plural agreement with diverses.
\smallskip
\item[\hypertarget{apparatus:FR-0206}{206}] Page \hyperlink{printpage:fr:208}{208}\quad\texttt{S45-SOL-0863}
\par\noindent\textbf{Imprimé:} un corps de nombre,
\par\noindent\textbf{Corrigé:} un corps de nombres,
\par\noindent Correct the mathematical term corps de nombres.
\smallskip
\item[\hypertarget{apparatus:FR-0207}{207}] Page \hyperlink{printpage:fr:208}{208}\quad\texttt{S45-SOL-0866}
\par\noindent\textbf{Imprimé:} soit en terme des orbites
\par\noindent\textbf{Corrigé:} soit en termes des orbites
\par\noindent Use the plural in the fixed expression en termes de.
\smallskip
\item[\hypertarget{apparatus:FR-0208}{208}] Page \hyperlink{printpage:fr:208}{208}\quad\texttt{S45-SOL-0867}
\par\noindent\textbf{Imprimé:} correspondant elle-même
\par\noindent\textbf{Corrigé:} correspondant elles-mêmes
\par\noindent The antecedent orbites is plural.
\smallskip
\item[\hypertarget{apparatus:FR-0209}{209}] Page \hyperlink{printpage:fr:208}{208}\quad\texttt{S45-SOL-0868}
\par\noindent\textbf{Imprimé:} celles telle que
\par\noindent\textbf{Corrigé:} celles telles que
\par\noindent Plural agreement with celles.
\smallskip
\item[\hypertarget{apparatus:FR-0210}{210}] Page \hyperlink{printpage:fr:209}{209}\quad\texttt{S45-SOL-0870}
\par\noindent\textbf{Imprimé:} $F^\ast$ dand $k^\ast$
\par\noindent\textbf{Corrigé:} $F^\ast$ dans $k^\ast$
\par\noindent Correct the misspelled preposition.
\smallskip
\item[\hypertarget{apparatus:FR-0211}{211}] Page \hyperlink{printpage:fr:210}{210}\quad\texttt{S45-SOL-0875}
\par\noindent\textbf{Imprimé:} des $x$ totalement positif
\par\noindent\textbf{Corrigé:} des $x$ totalement positifs
\par\noindent Plural agreement with the elements x.
\smallskip
\item[\hypertarget{apparatus:FR-0212}{212}] Page \hyperlink{printpage:fr:210}{210}\quad\texttt{S45-SOL-0877}
\par\noindent\textbf{Imprimé:} On pour toujours prendre
\par\noindent\textbf{Corrigé:} On peut toujours prendre
\par\noindent Correct the malformed verb.
\smallskip
\item[\hypertarget{apparatus:FR-0213}{213}] Page \hyperlink{printpage:fr:210}{210}\quad\texttt{S45-SOL-0879}
\par\noindent\textbf{Imprimé:} \(\chi_1|F_1=(N_{F_1/\dQ})^N\)
\par\noindent\textbf{Corrigé:} \(\chi_1|F_1'=(N_{F_1'/\dQ})^N\)
\par\noindent The restriction is to the totally real subfield, not to the entire CM field.
\smallskip
\item[\hypertarget{apparatus:FR-0214}{214}] Page \hyperlink{printpage:fr:210}{210}\quad\texttt{S45-SOL-0881}
\par\noindent\textbf{Imprimé:} idéaux fractionnaire
\par\noindent\textbf{Corrigé:} idéaux fractionnaires
\par\noindent Plural agreement.
\smallskip
\item[\hypertarget{apparatus:FR-0215}{215}] Page \hyperlink{printpage:fr:210}{210}\quad\texttt{S45-SOL-0882}
\par\noindent\textbf{Imprimé:} $k$ est sous torsion
\par\noindent\textbf{Corrigé:} $k$ est sans torsion
\par\noindent The fractional-ideal group is torsion-free, as used by the argument.
\smallskip
\item[\hypertarget{apparatus:FR-0216}{216}] Page \hyperlink{printpage:fr:211}{211}\quad\texttt{S45-SOL-0883}
\par\noindent\textbf{Imprimé:} Soient $F$ de $k$
\par\noindent\textbf{Corrigé:} Soient $F$ et $k$
\par\noindent The two fields are joined by et.
\smallskip
\item[\hypertarget{apparatus:FR-0217}{217}] Page \hyperlink{printpage:fr:212}{212}\quad\texttt{S45-SOL-0884}
\par\noindent\textbf{Imprimé:} groupe (fini) $\dmu$ racines
\par\noindent\textbf{Corrigé:} groupe (fini) $\dmu$ des racines
\par\noindent Restore the missing des.
\smallskip
\item[\hypertarget{apparatus:FR-0218}{218}] Page \hyperlink{printpage:fr:213}{213}\quad\texttt{S45-SOL-0889}
\par\noindent\textbf{Imprimé:} K\"ummer
\par\noindent\textbf{Corrigé:} Kummer
\par\noindent Correct the proper name.
\smallskip
\item[\hypertarget{apparatus:FR-0219}{219}] Page \hyperlink{printpage:fr:214}{214}\quad\texttt{S45-SOL-0892}
\par\noindent\textbf{Imprimé:} la famille $\lambda$ est défini
\par\noindent\textbf{Corrigé:} la famille $\lambda$ est définie
\par\noindent Agreement with famille.
\smallskip
\item[\hypertarget{apparatus:FR-0220}{220}] Page \hyperlink{printpage:fr:214}{214}\quad\texttt{S45-SOL-0893}
\par\noindent\textbf{Imprimé:} l'ensemble des homomorphisme dans
\par\noindent\textbf{Corrigé:} l'ensemble des homomorphismes dans
\par\noindent Plural agreement after des.
\smallskip
\item[\hypertarget{apparatus:FR-0221}{221}] Page \hyperlink{printpage:fr:214}{214}\quad\texttt{S45-SOL-0898}
\par\noindent\textbf{Imprimé:} plongements de $e_\beta$ dans $f$
\par\noindent\textbf{Corrigé:} plongements de $e_\beta$ dans $\bar f$
\par\noindent The embeddings are into the algebraic closure of the residue field.
\smallskip
\item[\hypertarget{apparatus:FR-0222}{222}] Page \hyperlink{printpage:fr:215}{215}\quad\texttt{S45-SOL-0899}
\par\noindent\textbf{Imprimé:} \(\lambda_{d(\beta)}\)
\par\noindent\textbf{Corrigé:} \(\lambda_{\alpha(\beta)}\)
\par\noindent The character is indexed by the orbit alpha(beta).
\smallskip
\item[\hypertarget{apparatus:FR-0223}{223}] Page \hyperlink{printpage:fr:215}{215}\quad\texttt{S45-SOL-0903}
\par\noindent\textbf{Imprimé:} \((\chi_\beta\circ \sigma_i^{-1})_{i\in I}\)
\par\noindent\textbf{Corrigé:} \((\chi_\beta\circ \sigma_i^{-1})_{i\in\beta}\)
\par\noindent The component torus uses the embeddings in beta.
\smallskip
\item[\hypertarget{apparatus:FR-0224}{224}] Page \hyperlink{printpage:fr:215}{215}\quad\texttt{S45-SOL-0904}
\par\noindent\textbf{Imprimé:} \(\sK_n((_n\bar\lambda_i)_{i\in \beta}) = \sF(\chi)\)
\par\noindent\textbf{Corrigé:} \(\sK_n((_n\bar\lambda_i)_{i\in I}) = \sF(\chi)\)
\par\noindent The global torus uses all embeddings I.
\smallskip
\item[\hypertarget{apparatus:FR-0225}{225}] Page \hyperlink{printpage:fr:215}{215}\quad\texttt{S45-SOL-0906}
\par\noindent\textbf{Imprimé:} \((\chi_\beta)_{\beta\in B}\)
\par\noindent\textbf{Corrigé:} \((\chi_\beta)_{\beta\in D}\)
\par\noindent The orbit set is D, not the undefined B.
\smallskip
\item[\hypertarget{apparatus:FR-0226}{226}] Page \hyperlink{printpage:fr:216}{216}\quad\texttt{S45-SOL-0908}
\par\noindent\textbf{Imprimé:} ensemble de densité $0$
\par\noindent\textbf{Corrigé:} ensemble de densité $1$
\par\noindent Degree-one prime ideals have density one among prime ideals of the number field.
\smallskip
\item[\hypertarget{apparatus:FR-0227}{227}] Page \hyperlink{printpage:fr:216}{216}\quad\texttt{S45-SOL-0912}
\par\noindent\textbf{Imprimé:} des homomorphisme algébriques
\par\noindent\textbf{Corrigé:} des homomorphismes algébriques
\par\noindent Plural agreement.
\smallskip
\item[\hypertarget{apparatus:FR-0228}{228}] Page \hyperlink{printpage:fr:217}{217}\quad\texttt{S45-ASTRA-S-0037}
\par\noindent\textbf{Imprimé:} $a\in A_j\subset (\dZ/n)^\ast$
\par\noindent\textbf{Corrigé:} $a\in A_j\subset\dZ/n$
\par\noindent The orbit family is in all of Z/n, as defined in 6.9(a).
\smallskip
\item[\hypertarget{apparatus:FR-0229}{229}] Page \hyperlink{printpage:fr:217}{217}\quad\texttt{S45-SOL-0919}
\par\noindent\textbf{Imprimé:} $a$ transformé
\par\noindent\textbf{Corrigé:} $a$ transforme
\par\noindent The group element a is the subject; use the present finite verb.
\smallskip
\item[\hypertarget{apparatus:FR-0230}{230}] Page \hyperlink{printpage:fr:217}{217}\quad\texttt{S45-SOL-0920}
\par\noindent\textbf{Imprimé:} sont invariant par $H$
\par\noindent\textbf{Corrigé:} sont invariants par $H$
\par\noindent Plural agreement.
\smallskip
\item[\hypertarget{apparatus:FR-0231}{231}] Page \hyperlink{printpage:fr:219}{219}\quad\texttt{S45-SOL-0930}
\par\noindent\textbf{Imprimé:} les énoncés cohomologiques de reflet les
\par\noindent\textbf{Corrigé:} les énoncés cohomologiques que reflètent les
\par\noindent Restore the grammatical relation between the statements and the identities reflecting them.
\smallskip
\item[\hypertarget{apparatus:FR-0232}{232}] Page \hyperlink{printpage:fr:220}{220}\quad\texttt{S45-SOL-0933}
\par\noindent\textbf{Imprimé:} en terme de sommes de Gauss
\par\noindent\textbf{Corrigé:} en termes de sommes de Gauss
\par\noindent Correct the fixed expression en termes de.
\smallskip
\item[\hypertarget{apparatus:FR-0233}{233}] Page \hyperlink{printpage:fr:220}{220}\quad\texttt{S45-SOL-0934}
\par\noindent\textbf{Imprimé:} \(\kloos_{\dF_{q^n},a}\)
\par\noindent\textbf{Corrigé:} \(\kloos_{\dF_q^n,a}\)
\par\noindent The original n-variable sum uses the split étale algebra F\_q\textasciicircum{}n.
\smallskip
\item[\hypertarget{apparatus:FR-0234}{234}] Page \hyperlink{printpage:fr:222}{222}\quad\texttt{S45-SOL-0940}
\par\noindent\textbf{Imprimé:} sur l'ouvert dense $\dA^1$ où
\par\noindent\textbf{Corrigé:} sur l'ouvert dense de $\dA^1$ où
\par\noindent Supply the preposition locating the dense open subset.
\smallskip
\item[\hypertarget{apparatus:FR-0235}{235}] Page \hyperlink{printpage:fr:224}{224}\quad\texttt{S45-SOL-0945}
\par\noindent\textbf{Imprimé:} $\eR^{n-1} \pi_! \sF(\psi\sigma)$ être unipotente
\par\noindent\textbf{Corrigé:} $\eR^{n-1} \pi_! \sF(\psi\sigma)$ d'être unipotente
\par\noindent Supply the preposition required by impose ... de.
\smallskip
\item[\hypertarget{apparatus:FR-0236}{236}] Page \hyperlink{printpage:fr:226}{226}\quad\texttt{S45-SOL-0951}
\par\noindent\textbf{Imprimé:} valeurs propres de $F$ est donc $q^{\frac{1-n}{2}}$
\par\noindent\textbf{Corrigé:} valeurs propres de $F$ est donc $q^{\frac{n-1}{2}}$
\par\noindent Undo the normalization factor of F to recover its eigenvalue.
\smallskip
\item[\hypertarget{apparatus:FR-0237}{237}] Page \hyperlink{printpage:fr:227}{227}\quad\texttt{S45-SOL-0959}
\par\noindent\textbf{Imprimé:} fonctions sur $X$ invariante par $\tau$
\par\noindent\textbf{Corrigé:} fonctions sur $V_a^1\times V_a^1$ invariantes par $\tau$
\par\noindent Invariant functions are on the space before taking the quotient; also restore plural agreement.
\smallskip
\item[\hypertarget{apparatus:FR-0238}{238}] Page \hyperlink{printpage:fr:228}{228}\quad\texttt{S45-SOL-0966}
\par\noindent\textbf{Imprimé:} Le groupe symétrique $\mathfrak{S}_m$
\par\noindent\textbf{Corrigé:} Le groupe symétrique $\mathfrak{S}_n$
\par\noindent The group permutes n coordinates.
\smallskip
\item[\hypertarget{apparatus:FR-0239}{239}] Page \hyperlink{printpage:fr:228}{228}\quad\texttt{S45-SOL-0969}
\par\noindent\textbf{Imprimé:} Il s'agit donc sur
\par\noindent\textbf{Corrigé:} Il agit donc sur
\par\noindent The intended verb is agir on the cohomology group.
\smallskip
\item[\hypertarget{apparatus:FR-0240}{240}] Page \hyperlink{printpage:fr:228}{228}\quad\texttt{S45-SOL-0971}
\par\noindent\textbf{Imprimé:} pour valeur propres
\par\noindent\textbf{Corrigé:} pour valeurs propres
\par\noindent Plural agreement.
\smallskip
\item[\hypertarget{apparatus:FR-0241}{241}] Page \hyperlink{printpage:fr:229}{229}\quad\texttt{S45-SOL-0973}
\par\noindent\textbf{Imprimé:} \(\h_0^\bullet(\dG_a,\sF(\psi))\)
\par\noindent\textbf{Corrigé:} \(\h_c^\bullet(\dG_a,\sF(\psi))\)
\par\noindent Use compactly supported cohomology in the Künneth vanishing.
\smallskip
\item[\hypertarget{apparatus:FR-0242}{242}] Page \hyperlink{printpage:fr:230}{230}\quad\texttt{S45-SOL-0978}
\par\noindent\textbf{Imprimé:} groupe de Weil $W$
\par\noindent\textbf{Corrigé:} groupe de Weyl $W$
\par\noindent The group acting on the maximal torus is the Weyl group.
\smallskip
\item[\hypertarget{apparatus:FR-0243}{243}] Page \hyperlink{printpage:fr:232}{232}\quad\texttt{S45-ASTRA-S-0049}
\par\noindent\textbf{Imprimé:} Weil
\par\noindent\textbf{Corrigé:} Weyl
\par\noindent Correct the proper name of the reflection groups.
\smallskip
\item[\hypertarget{apparatus:FR-0244}{244}] Page \hyperlink{printpage:fr:234}{234}\quad\texttt{S45-SOL-0981}
\par\noindent\textbf{Imprimé:} \(E_1^{p q}\)
\par\noindent\textbf{Corrigé:} \(E_2^{p q}\)
\par\noindent The first display on p.234 visibly has E\_1. The displayed derived-functor cohomology terms constitute the E\_2 page, as in the Leray display immediately below; retain 1 diplomatically and correct to 2.
\smallskip
\item[\hypertarget{apparatus:FR-0245}{245}] Page \hyperlink{printpage:fr:234}{234}\quad\texttt{S45-SOL-0982}
\par\noindent\textbf{Imprimé:} dérivées à l'avantage
\par\noindent\textbf{Corrigé:} dérivées a l'avantage
\par\noindent The first prose paragraph visibly prints à l’avantage; the subject infinitive requires the finite verb a. This French accent error does not change the English.
\smallskip
\item[\hypertarget{apparatus:FR-0246}{246}] Page \hyperlink{printpage:fr:234}{234}\quad\texttt{S45-SOL-0989}
\par\noindent\textbf{Imprimé:} \(\eR g^! K(n)[2 n]\)
\par\noindent\textbf{Corrigé:} \(\eR g^! K \iso g^\ast K(n)[2 n]\)
\par\noindent The bottom line contains Rg\textasciicircum{}!K(n)[2n] without an equality or pullback. Smooth duality and the stated relative dimension require Rg\textasciicircum{}!K ≃ g\textasciicircum{}*K(n)[2n].
\smallskip
\item[\hypertarget{apparatus:FR-0247}{247}] Page \hyperlink{printpage:fr:235}{235}\quad\texttt{S45-SOL-0990}
\par\noindent\textbf{Imprimé:} \(K\to \eR j_\ast j^\ast K\)
\par\noindent\textbf{Corrigé:} \(L\to \eR j_\ast j^\ast L\)
\par\noindent At the top of p.235 the variable L is declared, but the displayed arrow contains K twice. The natural map for the same object L requires L in both positions.
\smallskip
\item[\hypertarget{apparatus:FR-0248}{248}] Page \hyperlink{printpage:fr:235}{235}\quad\texttt{S45-SOL-0992}
\par\noindent\textbf{Imprimé:} \(\eR j_\ast \Hom(\sF_1,\eR j^! \sG)\)
\par\noindent\textbf{Corrigé:} \(\eR j_\ast \rHom(\sF_1,\eR j^! \sG)\)
\par\noindent The right side of the printed derived Hom identity lacks R before internal Hom. Its left side and the subsequent derived assertion require the derived internal Hom; the correction is substantive, not a choice of sheaf typography.
\smallskip
\item[\hypertarget{apparatus:FR-0249}{249}] Page \hyperlink{printpage:fr:236}{236}\quad\texttt{S45-SOL-0996}
\par\noindent\textbf{Imprimé:} Soit $F:X\to Y$ un morphisme de $S$-schémas de type fini
\par\noindent\textbf{Corrigé:} Soit $f:X\to Y$ un morphisme de $S$-schémas de type fini
\par\noindent Theorem1.9 starts with capital F, while both its conclusions use lowercase f. The capital is a misprint; preserve it only diplomatically.
\smallskip
\item[\hypertarget{apparatus:FR-0250}{250}] Page \hyperlink{printpage:fr:239}{239}\quad\texttt{S45-SOL-0998}
\par\noindent\textbf{Imprimé:} telle 
    que l'image réciproque
\par\noindent\textbf{Corrigé:} tel 
    que l'image réciproque
\par\noindent The opening item(c) visibly reads telle que, agreeing with no feminine antecedent: revêtement is masculine. Correct tel in French only.
\smallskip
\item[\hypertarget{apparatus:FR-0251}{251}] Page \hyperlink{printpage:fr:239}{239}\quad\texttt{S45-SOL-0999}
\par\noindent\textbf{Imprimé:} \(Y\subset \dA_S^n\)
\par\noindent\textbf{Corrigé:} \(Y\subset \dA_S^N\)
\par\noindent At the bottom of239 the affine embedding uses n, already fixed as the relative-dimension bound in(*\_n). An arbitrary affine scheme of dimension at most n need not embed in affine n-space; distinguish the embedding dimension N.
\smallskip
\item[\hypertarget{apparatus:FR-0252}{252}] Page \hyperlink{printpage:fr:240}{240}\quad\texttt{S45-SOL-1001}
\par\noindent\textbf{Imprimé:} elle valent
\par\noindent\textbf{Corrigé:} elles valent
\par\noindent The next line visibly has elle valent, referring to les conclusions. Restore plural elles only in corrected French.
\smallskip
\item[\hypertarget{apparatus:FR-0253}{253}] Page \hyperlink{printpage:fr:240}{240}\quad\texttt{S45-SOL-1004}
\par\noindent\textbf{Imprimé:} changement de Base en $S$, et de même
\par\noindent\textbf{Corrigé:} changement de base en $S$, et de même
\par\noindent The final line visibly capitalizes Base within changement de Base. This isolated capital is a typographic misprint, retained in diplomatic French and lowered in corrected French.
\smallskip
\item[\hypertarget{apparatus:FR-0254}{254}] Page \hyperlink{printpage:fr:243}{243}\quad\texttt{S45-SOL-1015}
\par\noindent\textbf{Imprimé:} \ref{VII:2-13} est toujours égal à $S$.
\par\noindent\textbf{Corrigé:} \ref{VII:2-13} est toujours égal à $S$).
\par\noindent The last sentence opens a parenthesis after limite and never closes it before the final full stop. Supply the missing parenthesis only in corrected forms.
\smallskip
\item[\hypertarget{apparatus:FR-0255}{255}] Page \hyperlink{printpage:fr:245}{245}\quad\texttt{S45-SOL-1023}
\par\noindent\textbf{Imprimé:} l'une des projection 
\par\noindent\textbf{Corrigé:} l'une des projections 
\par\noindent The paragraph beginning Prouvons3.3 visibly has l’une des projection. The plural noun projections is required after l’une des. Soient governs f, X′ and F′ and remains plural.
\smallskip
\item[\hypertarget{apparatus:FR-0256}{256}] Page \hyperlink{printpage:fr:246}{246}\quad\texttt{S45-SOL-1031}
\par\noindent\textbf{Imprimé:} \(\h^p(X_s,\eR^i\Psi_\eta(\sF))\)
\par\noindent\textbf{Corrigé:} \(\h^p(X_s,\eR^q\Psi_\eta(\sF))\)
\par\noindent The bottom spectral sequence prints i as the nearby-cycle degree, although the term is E\_2\textasciicircum{}\{pq\} and abuts in degree p+q. The nearby-cycle degree must be q; subsequent G\textasciicircum{}q and H\textasciicircum{}q confirm it.
\smallskip
\item[\hypertarget{apparatus:FR-0257}{257}] Page \hyperlink{printpage:fr:248}{248}\quad\texttt{S45-SOL-1035}
\par\noindent\textbf{Imprimé:} $\eR\Gamma(X_s,\Delta)=0$
\par\noindent\textbf{Corrigé:} $\eR\Gamma(X_{s'}',\Delta)=0$
\par\noindent Assertion(B) visibly uses X\_s. Its Δ was defined from the base-change morphism on X′\_\{s′\}; paragraph3.9 itself computes RΓ(X′\_\{s′\},Δ). Correct the space accordingly.
\smallskip
\item[\hypertarget{apparatus:FR-0258}{258}] Page \hyperlink{printpage:fr:248}{248}\quad\texttt{S45-SOL-1042}
\par\noindent\textbf{Imprimé:} les $\eR^i f_\ast \sF$ sont constructible au-dessus
\par\noindent\textbf{Corrigé:} les $\eR^i f_\ast \sF$ sont constructibles au-dessus
\par\noindent In3.10 the printed predicate after les R\textasciicircum{}i f\_*F is singular constructible. Correct plural agreement in French; the released English has no corresponding number defect.
\smallskip
\item[\hypertarget{apparatus:FR-0259}{259}] Page \hyperlink{printpage:fr:248}{248}\quad\texttt{S45-SOL-1043}
\par\noindent\textbf{Imprimé:} leur restriction 
aux $Y_t$
\par\noindent\textbf{Corrigé:} leurs restrictions 
aux $Y_t$
\par\noindent The next sentence visibly has leur restriction … sont constructibles. The plural verb and the family of fibers require leurs restrictions; this does not alter the English meaning.
\smallskip
\item[\hypertarget{apparatus:FR-0260}{260}] Page \hyperlink{printpage:fr:250}{250}\quad\texttt{S45-SOL-1045}
\par\noindent\textbf{Imprimé:} modules à droites
\par\noindent\textbf{Corrigé:} modules à droite
\par\noindent Paragraph4.1 explicitly prints à droites. The locution for right modules is à droite; remove only the spurious plural s in corrected French.
\smallskip
\item[\hypertarget{apparatus:FR-0261}{261}] Page \hyperlink{printpage:fr:250}{250}\quad\texttt{S45-SOL-1047}
\par\noindent\textbf{Imprimé:} $K\in\ob\eD_\text{ctf}^b(S,A)$, on pose
\par\noindent\textbf{Corrigé:} $K\in\ob\eD_\text{ctf}^b(X,A)$, on pose
\par\noindent Paragraph4.2 assigns K to D\_ctf\textasciicircum{}b(S,A), although K is dualized against K\_X and theorem4.3 concerns complexes on X. The printed S is a base-space error; use X in corrected forms.
\smallskip
\item[\hypertarget{apparatus:FR-0262}{262}] Page \hyperlink{printpage:fr:253}{253}\quad\texttt{S45-SOL-1056}
\par\noindent\textbf{Imprimé:} \(X_\bullet=X\times_S \bar S_\bullet\)
\par\noindent\textbf{Corrigé:} \(\bar X_\bullet=X\times_S \bar S_\bullet\)
\par\noindent The printed definition uses unbarred X\_bullet, but the very next projection is from barred X\_bullet. Supply the bar at the definition; retain the printed fiber product X×\_S barred S\_bullet, which is already correct.
\smallskip
\item[\hypertarget{apparatus:FR-0263}{263}] Page \hyperlink{printpage:fr:253}{253}\quad\texttt{S45-ASTRA-V-0004}
\par\noindent\textbf{Imprimé:} Posons $G|X'=G'$. Le bicomplexe 
${g_\bullet}_\ast g_\bullet^\ast G'$, de colonnes les 
${g_n}_\ast g_n^\ast G'$ est une résolution de $G'$
\par\noindent\textbf{Corrigé:} Posons $E|X'=E'$. Le bicomplexe 
${g_\bullet}_\ast g_\bullet^\ast E'$, de colonnes les 
${g_n}_\ast g_n^\ast E'$ est une résolution de $E'$
\par\noindent The proposition fixes E and constructs the base-change cone from E. No object G is introduced before the resolution, which must resolve E′ to compute that cone. The scan clearly prints G in all five places.
\smallskip
\item[\hypertarget{apparatus:FR-0264}{264}] Page \hyperlink{printpage:fr:255}{255}\quad\texttt{S45-SOL-1062}
\par\noindent\textbf{Imprimé:} $\widetilde D$ de $D$ dans $Y$
\par\noindent\textbf{Corrigé:} $\widetilde D$ de $D$ dans $\widetilde Y$
\par\noindent The proof prints the inverse image D-tilde in Y, although the newly constructed covering is Y-tilde→Y. The inverse-image divisor belongs on Y-tilde.
\smallskip
\item[\hypertarget{apparatus:FR-0265}{265}] Page \hyperlink{printpage:fr:255}{255}\quad\texttt{S45-SOL-1063}
\par\noindent\textbf{Imprimé:} \(f_i:Y\setminus D_i \to D\)
\par\noindent\textbf{Corrigé:} \(f_i:Y\setminus D_i \to Y\)
\par\noindent The printed inclusion f\_i:Y−D\_i→D has target D. Its domain is an open subset of Y and the stated canonical inclusion must end at Y; the following direct-image formulas also live on Y.
\smallskip
\item[\hypertarget{apparatus:FR-0266}{266}] Page \hyperlink{printpage:fr:257}{257}\quad\texttt{S45-SOL-1070}
\par\noindent\textbf{Imprimé:} \(\eR\Gamma(\bar X,E|\bar X) (=\eR\bar f_\ast (E|\bar X)_s \to\)
\par\noindent\textbf{Corrigé:} \(\eR\Gamma(\bar X,E|\bar X) (=\eR\bar f_\ast (E|\bar X)_s) \to\)
\par\noindent The first equation on257 visibly opens a parenthesis before the equality and leaves it unclosed. Miller removes that opening mark. Restore it diplomatically; close the parenthetical equality before the arrow in corrected forms.
\smallskip
\item[\hypertarget{apparatus:FR-0267}{267}] Page \hyperlink{printpage:fr:257}{257}\quad\texttt{S45-SOL-1076}
\par\noindent\textbf{Imprimé:} \(\spe(u')\spe(u'')\)
\par\noindent\textbf{Corrigé:} \(\spe(u')\spe(u)\)
\par\noindent The specialization formula prints sp(u′)sp(u″), although the arrows just defined are u′ and u only. Functorial contravariant composition requires sp(u′)sp(u).
\smallskip
\item[\hypertarget{apparatus:FR-0268}{268}] Page \hyperlink{printpage:fr:259}{259}\quad\texttt{S45-SOL-1085}
\par\noindent\textbf{Imprimé:} \(\eR\Gamma(Y(y)_t,(\eR {f_{(x)}}_\ast (E|X(x))|Y(y)_t)\)
\par\noindent\textbf{Corrigé:} \(\eR\Gamma\left(Y(y)_t,\left(\eR {f_{(x)}}_\ast (E|X(x))\right)|Y(y)_t\right)\)
\par\noindent The right-hand cohomology term visibly opens a parenthesis before Rf\_(x)* but does not close that grouping before restriction. Supply the missing closing parenthesis around the complex in corrected forms; preserve the printed defect diplomatically.
\smallskip
\item[\hypertarget{apparatus:FR-0269}{269}] Page \hyperlink{printpage:fr:260}{260}\quad\texttt{S45-SOL-1086}
\par\noindent\textbf{Imprimé:} \(\eR f_\ast'(E X')\)
\par\noindent\textbf{Corrigé:} \(\eR f_\ast'(E|X')\)
\par\noindent The first proof paragraph prints Rf′\_*(E X′), omitting the restriction bar between E and X′. The argument and surrounding notation require E|X′; Miller has silently supplied it.
\smallskip
\item[\hypertarget{apparatus:FR-0270}{270}] Page \hyperlink{printpage:fr:260}{260}\quad\texttt{S45-SOL-1087}
\par\noindent\textbf{Imprimé:} image de $s$ 
(resp.\ $t$) dans $S$
\par\noindent\textbf{Corrigé:} image de $s'$ 
(resp.\ $t'$) dans $S$
\par\noindent The scan calls s and t the images of s and t, without primes. The source points just defined lie on S′ and are s′,t′; the missing primes must be supplied in the corrected text.
\smallskip
\item[\hypertarget{apparatus:FR-0271}{271}] Page \hyperlink{printpage:fr:260}{260}\quad\texttt{S45-SOL-1089}
\par\noindent\textbf{Imprimé:} \(\eR\Gamma(X(x)_t,(E|X(x)_t)\lotimes_A M\)
\par\noindent\textbf{Corrigé:} \(\eR\Gamma(X(x)_t,(E|X(x)_t)\lotimes_A M)\)
\par\noindent The final restriction formula on260 has an opening RΓ parenthesis with no terminal closing parenthesis after M. Miller silently closes it; retain the missing mark diplomatically and keep the closing mark in corrected forms.
\smallskip
\item[\hypertarget{apparatus:FR-0272}{272}] Page \hyperlink{printpage:fr:264}{264}\quad\texttt{S45-SOL-1096}
\par\noindent\textbf{Imprimé:} \(\beta\circ \alpha = \beta\circ T^i(\alpha)\)
\par\noindent\textbf{Corrigé:} \(\beta\circ \alpha = T^i(\beta)\circ\alpha\)
\par\noindent The scan visibly prints β∘T\textasciicircum{}i(α). With Hom\textasciicircum{}i(X,Y)=Hom(X,T\textasciicircum{}iY), the composable ordinary arrows are α followed by T\textasciicircum{}i(β), so correct to T\textasciicircum{}i(β)∘α.
\smallskip
\item[\hypertarget{apparatus:FR-0273}{273}] Page \hyperlink{printpage:fr:265}{265}\quad\texttt{S45-ASTRA-V-0019}
\par\noindent\textbf{Imprimé:} \(F\circ T = T'\circ F\)
\par\noindent\textbf{Corrigé:} \(F\circ T \simeq T'\circ F\)
\par\noindent The1976 footnote explicitly says that the original equality requirement was changed to an isomorphism. The unchanged equality sign in the displayed specification should be an isomorphism sign; the Codex proposal is accepted.
\smallskip
\item[\hypertarget{apparatus:FR-0274}{274}] Page \hyperlink{printpage:fr:265}{265}\quad\texttt{S45-SOL-1100}
\par\noindent\textbf{Imprimé:} $u:X\leftrightarrow Y$, $v:Y\to Z$
\par\noindent\textbf{Corrigé:} $u:X\to Y$, $v:Y\to Z$
\par\noindent The arrow in the printed list for u has arrowheads at both ends. The same page’s triangle and all ensuing axioms require the single morphism u:X→Y; retain the two-headed misprint diplomatically.
\smallskip
\item[\hypertarget{apparatus:FR-0275}{275}] Page \hyperlink{printpage:fr:267}{267}\quad\texttt{S45-SOL-1108}
\par\noindent\textbf{Imprimé:} {\footnotesize\ttfamily\raggedright \textbackslash{}\allowbreak{}hom(\allowbreak{}X,\allowbreak{}M)\allowbreak{} \textbackslash{}\allowbreak{}ar[\allowbreak{}r]\allowbreak{}\textasciicircum{}\allowbreak{}-\allowbreak{}\{\allowbreak{}\textbackslash{}\allowbreak{}hom(\allowbreak{}w,\allowbreak{}M)\allowbreak{}\}\allowbreak{}}
\par\noindent\textbf{Corrigé:} {\footnotesize\ttfamily\raggedright \textbackslash{}\allowbreak{}hom\textasciicircum{}\allowbreak{}i(\allowbreak{}X,\allowbreak{}M)\allowbreak{} \textbackslash{}\allowbreak{}ar[\allowbreak{}r]\allowbreak{}\textasciicircum{}\allowbreak{}-\allowbreak{}\{\allowbreak{}\textbackslash{}\allowbreak{}hom(\allowbreak{}w,\allowbreak{}M)\allowbreak{}\}\allowbreak{}}
\par\noindent The third term in the second exact sequence has Hom(X,M) without its degree superscript. The neighboring terms and degree-zero arrows show that the omitted superscript is i.
\smallskip
\item[\hypertarget{apparatus:FR-0276}{276}] Page \hyperlink{printpage:fr:268}{268}\quad\texttt{S45-SOL-1112}
\par\noindent\textbf{Imprimé:} Les morphismes de $\cA$ seront
\par\noindent\textbf{Corrigé:} Les morphismes de $\eC(\cA)$ seront
\par\noindent The second bullet on268 prints A alone. Morphisms of complexes are the morphisms of C(A), as both the immediately preceding definition and all subsequent constructions require.
\smallskip
\item[\hypertarget{apparatus:FR-0277}{277}] Page \hyperlink{printpage:fr:269}{269}\quad\texttt{S45-SOL-1113}
\par\noindent\textbf{Imprimé:} automorphisme additif de $\cA$
\par\noindent\textbf{Corrigé:} automorphisme additif de $\eC(\cA)$
\par\noindent The first prose paragraph on269 again has A alone. T is the translation automorphism on C(A), not on the original additive category A.
\smallskip
\item[\hypertarget{apparatus:FR-0278}{278}] Page \hyperlink{printpage:fr:270}{270}\quad\texttt{S45-SOL-1118}
\par\noindent\textbf{Imprimé:} couple d'objet
\par\noindent\textbf{Corrigé:} couple d'objets
\par\noindent The scan has singular objet after couple. French grammar requires objets; the English phrase pair of objects already conveys the intended plural.
\smallskip
\item[\hypertarget{apparatus:FR-0279}{279}] Page \hyperlink{printpage:fr:271}{271}\quad\texttt{S45-SOL-1120}
\par\noindent\textbf{Imprimé:} engendrées par ses familles
\par\noindent\textbf{Corrigé:} engendrées par ces familles
\par\noindent The first bullet has ses familles. It points back to the three just-listed families of bounded complexes; the demonstrative ces is the intended reading.
\smallskip
\item[\hypertarget{apparatus:FR-0280}{280}] Page \hyperlink{printpage:fr:271}{271}\quad\texttt{S45-SOL-1121}
\par\noindent\textbf{Imprimé:} On paiera aussi
\par\noindent\textbf{Corrigé:} On parlera aussi
\par\noindent The scan unambiguously reads On paiera aussi. The paragraph introduces more subcategories rather than payment; parlera is the intended verb.
\smallskip
\item[\hypertarget{apparatus:FR-0281}{281}] Page \hyperlink{printpage:fr:273}{273}\quad\texttt{S45-SOL-1124}
\par\noindent\textbf{Imprimé:} {\footnotesize\ttfamily\raggedright 0 \& d\_\allowbreak{}\{\allowbreak{}Y\textasciicircum{}\allowbreak{}\textbackslash{}\allowbreak{}bullet\}\allowbreak{}}
\par\noindent\textbf{Corrigé:} {\footnotesize\ttfamily\raggedright 0 \& d\_\allowbreak{}\{\allowbreak{}Z\textasciicircum{}\allowbreak{}\textbackslash{}\allowbreak{}bullet\}\allowbreak{}}
\par\noindent The lower-right matrix entry visibly has Y. The decomposition Y=X+Z requires d\_Z in that corner. Preserve the printed Y as a mathematical misprint in diplomatic forms.
\smallskip
\item[\hypertarget{apparatus:FR-0282}{282}] Page \hyperlink{printpage:fr:274}{274}\quad\texttt{S45-ASTRA-V-0021}
\par\noindent\textbf{Imprimé:} \(F^{i+-}(X)\)
\par\noindent\textbf{Corrigé:} \(F^{i+1}(X)\)
\par\noindent A high-resolution scan crop shows a horizontal stroke after the plus, not a1. The exact sequence shifts cohomological degree by one. Accept the Codex correction; preserve the defective printed exponent diplomatically.
\smallskip
\item[\hypertarget{apparatus:FR-0283}{283}] Page \hyperlink{printpage:fr:275}{275}\quad\texttt{S45-SOL-1126}
\par\noindent\textbf{Imprimé:} \(P^\bullet(Z^\bullet,f)\)
\par\noindent\textbf{Corrigé:} \(P^\ast(Z^\bullet,f)\)
\par\noindent The parenthetical functor on complexes visibly has a bullet, although P\textasciicircum{}bullet is not defined until the ensuing quotient construction. At this point the functor is P\textasciicircum{}ast; the English already has the corrected symbol.
\smallskip
\item[\hypertarget{apparatus:FR-0284}{284}] Page \hyperlink{printpage:fr:275}{275}\quad\texttt{S45-ASTRA-V-0039}
\par\noindent\textbf{Imprimé:} \(\hom^\circ:\eK(\cA)^\circ\)
\par\noindent\textbf{Corrigé:} \(\hom^\bullet:\eK(\cA)^\circ\)
\par\noindent The enlarged275 crop shows the same hollow circle as the adjacent opposite-category superscript. Page309 explicitly refers back to this construction and writes Hom with a filled bullet, consistently with the complex-valued functor and its arguments. The distinct printed glyph is retained diplomatically rather than silently normalized.
\smallskip
\item[\hypertarget{apparatus:FR-0285}{285}] Page \hyperlink{printpage:fr:280}{280}\quad\texttt{S45-SOL-1138}
\par\noindent\textbf{Imprimé:} transformant morphisme de $S$
\par\noindent\textbf{Corrigé:} transformant tout morphisme de $S$
\par\noindent The printed universal-property sentence omits a determiner before morphisme. Add tout in the corrected French forms; English takes morphisms to isomorphisms already expresses the universal condition.
\smallskip
\item[\hypertarget{apparatus:FR-0286}{286}] Page \hyperlink{printpage:fr:281}{281}\quad\texttt{S45-SOL-1145}
\par\noindent\textbf{Imprimé:} un objet de $S_Y$, $\underline b$
\par\noindent\textbf{Corrigé:} $t$ un objet de $S_Y$, $\underline b$
\par\noindent The scan has et un objet de S\_Y, without naming that object. The subsequent source(t) and diagram show that the omitted variable is t; insert it only in corrected forms.
\smallskip
\item[\hypertarget{apparatus:FR-0287}{287}] Page \hyperlink{printpage:fr:282}{282}\quad\texttt{S45-ASTRA-V-0025}
\par\noindent\textbf{Imprimé:} ne sont autre que
\par\noindent\textbf{Corrigé:} ne sont autres que
\par\noindent The scan has singular autre governed by plural triangles. Correct the agreement only after restoring the actual printed initial letter. The English has no corresponding agreement defect.
\smallskip
\item[\hypertarget{apparatus:FR-0288}{288}] Page \hyperlink{printpage:fr:286}{286}\quad\texttt{S45-ASTRA-V-0026}
\par\noindent\textbf{Imprimé:} où $f=i\circ h$ et où $g\circ i=0$.
\par\noindent\textbf{Corrigé:} où $N'$ est un objet de $\cB$, $f=i\circ h$ et $g\circ i=0$.
\par\noindent The scan omits N′∈B. Without it the condition is vacuous: choose N′=X, h=id\_X and i=f. The equivalence with surjectivity therefore requires this omitted membership hypothesis. Accept the later corrected-tree proposal.
\smallskip
\item[\hypertarget{apparatus:FR-0289}{289}] Page \hyperlink{printpage:fr:288}{288}\quad\texttt{S45-ASTRA-V-0027}
\par\noindent\textbf{Imprimé:} La catégorie $\cB^\bot/X$
\par\noindent\textbf{Corrigé:} La catégorie $X\backslash\cB^\bot$
\par\noindent The printed slash category repeats the notation for objects over X, whereas the same clause explicitly says under X and asks for an initial object. The corrected-tree replacement X\textbackslash{}B⊥ gives the comma category described by the words and the adjunction. Preserve the printed slash in diplomatic forms.
\smallskip
\item[\hypertarget{apparatus:FR-0290}{290}] Page \hyperlink{printpage:fr:290}{290}\quad\texttt{S45-SOL-1169}
\par\noindent\textbf{Imprimé:} \(Q^\ast\circ Q\iso I^\bot\circ ^\ast i^\bot\)
\par\noindent\textbf{Corrigé:} \(Q^\ast\circ Q\iso i^\bot\circ ^\ast i^\bot\)
\par\noindent The first textual isomorphism on290 visibly uses capital I\textasciicircum{}⊥. The inclusion functor was defined as lowercase i\textasciicircum{}⊥ on289, and the later triangle uses that lowercase symbol, so this is a print misprint.
\smallskip
\item[\hypertarget{apparatus:FR-0291}{291}] Page \hyperlink{printpage:fr:290}{290}\quad\texttt{S45-SOL-1170}
\par\noindent\textbf{Imprimé:} {\footnotesize\ttfamily\raggedright \& \& I\textasciicircum{}\allowbreak{}\textbackslash{}\allowbreak{}bot\textbackslash{}\allowbreak{}circ \textasciicircum{}\allowbreak{}\textbackslash{}\allowbreak{}ast i\textasciicircum{}\allowbreak{}\textbackslash{}\allowbreak{}bot}
\par\noindent\textbf{Corrigé:} {\footnotesize\ttfamily\raggedright \& \& i\textasciicircum{}\allowbreak{}\textbackslash{}\allowbreak{}bot\textbackslash{}\allowbreak{}circ \textasciicircum{}\allowbreak{}\textbackslash{}\allowbreak{}ast i\textasciicircum{}\allowbreak{}\textbackslash{}\allowbreak{}bot}
\par\noindent The lower-right vertex of the first diagram repeats the same capital I\textasciicircum{}⊥ misprint. Correct it consistently with the functor i\textasciicircum{}⊥.
\smallskip
\item[\hypertarget{apparatus:FR-0292}{292}] Page \hyperlink{printpage:fr:290}{290}\quad\texttt{S45-SOL-1172}
\par\noindent\textbf{Imprimé:} \(\delta(T X_0) = - T \delta(X)\)
\par\noindent\textbf{Corrigé:} \(\delta(T X) = - T \delta(X)\)
\par\noindent The final identity visibly adds a small round subscript to the left X. There is no distinguished object X\_0, and the right side is evaluated at arbitrary X. Remove the subscript in corrected forms while retaining the printed final δ notation.
\smallskip
\item[\hypertarget{apparatus:FR-0293}{293}] Page \hyperlink{printpage:fr:292}{292}\quad\texttt{S45-SOL-1177}
\par\noindent\textbf{Imprimé:} \(\eK^{b,b}\)
\par\noindent\textbf{Corrigé:} \(\eK^{b,b}(\cA)\)
\par\noindent The first natural-functor formula on292 omits(A) after K\textasciicircum{}\{b,b\}; the second formula and the definition on291 include it. Restore the argument only in corrected forms.
\smallskip
\item[\hypertarget{apparatus:FR-0294}{294}] Page \hyperlink{printpage:fr:293}{293}\quad\texttt{S45-SOL-1178}
\par\noindent\textbf{Imprimé:} sera encore noté $D$
\par\noindent\textbf{Corrigé:} sera encore notée $D$
\par\noindent The first sentence on293 has noté, but the subject is the feminine restriction. Correct the agreement in French; English denoted is unaffected.
\smallskip
\item[\hypertarget{apparatus:FR-0295}{295}] Page \hyperlink{printpage:fr:293}{293}\quad\texttt{S45-SOL-1180}
\par\noindent\textbf{Imprimé:} $m\geqslant 0$
\par\noindent\textbf{Corrigé:} $m>0$
\par\noindent The print includes m=0 in the vanishing claim. At m=0 the identity of a nonzero object survives under the fully faithful embedding stated immediately above, so the correct range is m>0.
\smallskip
\item[\hypertarget{apparatus:FR-0296}{296}] Page \hyperlink{printpage:fr:295}{295}\quad\texttt{S45-SOL-1186}
\par\noindent\textbf{Imprimé:} un morphisme $f:X^\bullet \to Y^\bullet$ et un
\par\noindent\textbf{Corrigé:} un morphisme $f:X^\bullet \to Y^\bullet$ est un
\par\noindent The scan prints et un isomorphisme. The criterion is the assertion that f is an isomorphism, so est is required; preserve the conjunction in diplomatic French and English.
\smallskip
\item[\hypertarget{apparatus:FR-0297}{297}] Page \hyperlink{printpage:fr:296}{296}\quad\texttt{S45-SOL-1190}
\par\noindent\textbf{Imprimé:} ambiguité
\par\noindent\textbf{Corrigé:} ambiguïté
\par\noindent The print spells ambiguité without the required diaeresis on ï. Correct the French spelling; English ambiguity is unchanged.
\smallskip
\item[\hypertarget{apparatus:FR-0298}{298}] Page \hyperlink{printpage:fr:296}{296}\quad\texttt{S45-SOL-1191}
\par\noindent\textbf{Imprimé:} d'objets de $\cA$ ou pour
\par\noindent\textbf{Corrigé:} d'objets de $\cA$ où pour
\par\noindent The print has ou without the grave accent. This is the relative adverb describing the sequence, so où is required, not the alternative conjunction.
\smallskip
\item[\hypertarget{apparatus:FR-0299}{299}] Page \hyperlink{printpage:fr:298}{298}\quad\texttt{S45-SOL-1194}
\par\noindent\textbf{Imprimé:} \(\ext^0(X^\bullet,Y^\bullet) \iso \varinjlim_{Y^\bullet/\mathsf{Qis}}\)
\par\noindent\textbf{Corrigé:} \(\ext^0(X^\bullet,Y^\bullet) \iso \varinjlim_{Y^\bullet/\mathsf{Qis}^+}\)
\par\noindent In the bounded-below assertion the printed limit is indexed by all Qis. The companion bounded-above formula uses Qis\textasciicircum{}-; the corresponding strengthened bounded-below formula requires Qis\textasciicircum{}+. Retain the omitted plus diplomatically and restore it in corrected forms.
\smallskip
\item[\hypertarget{apparatus:FR-0300}{300}] Page \hyperlink{printpage:fr:298}{298}\quad\texttt{S45-SOL-1197}
\par\noindent\textbf{Imprimé:} suffisamment d'injectif et
\par\noindent\textbf{Corrigé:} suffisamment d'injectifs et
\par\noindent The print uses singular injectif after suffisamment de. Correct to injectifs in French; the English already uses the ordinary plural.
\smallskip
\item[\hypertarget{apparatus:FR-0301}{301}] Page \hyperlink{printpage:fr:302}{302}\quad\texttt{S45-SOL-1207}
\par\noindent\textbf{Imprimé:} \(\h{}\circ \eK f\)
\par\noindent\textbf{Corrigé:} \(\h^0{}\circ \eK f\)
\par\noindent The unbounded cohomological definition prints H∘Kf with no exponent. It is the zero-cohomology functor, as in the immediately preceding bounded definition; restore the omitted0 in corrected forms.
\smallskip
\item[\hypertarget{apparatus:FR-0302}{302}] Page \hyperlink{printpage:fr:303}{303}\quad\texttt{S45-SOL-1213}
\par\noindent\textbf{Imprimé:} entiers rationnel soit
\par\noindent\textbf{Corrigé:} entiers rationnels soit
\par\noindent The scan prints entiers rationnel, with no plural s on the adjective. Correct this French agreement; rational integers is already plural in English.
\smallskip
\item[\hypertarget{apparatus:FR-0303}{303}] Page \hyperlink{printpage:fr:304}{304}\quad\texttt{S45-SOL-1216}
\par\noindent\textbf{Imprimé:} ensembles d'objet sont
\par\noindent\textbf{Corrigé:} ensembles d'objets sont
\par\noindent Remarks2.2 prints ensembles d’objet. The sets consist of objects, so the French noun takes the plural; English sets of objects is already correct.
\smallskip
\item[\hypertarget{apparatus:FR-0304}{304}] Page \hyperlink{printpage:fr:305}{305}\quad\texttt{S45-SOL-1221}
\par\noindent\textbf{Imprimé:} $Q\circ \Phi(Y)$
\par\noindent\textbf{Corrigé:} $Q'\circ \Phi(Y)$
\par\noindent Item(c) prints Q∘Φ(Y), but Φ(Y) belongs to K(B) and must be sent to D(B) by Q′. Q has source K(A). The corrected-tree and English reading Q′ is justified by these explicitly given domains.
\smallskip
\item[\hypertarget{apparatus:FR-0305}{305}] Page \hyperlink{printpage:fr:305}{305}\quad\texttt{S45-SOL-1222}
\par\noindent\textbf{Imprimé:} \(\varinjlim_{X/\mathsf{Qis}^\ast} Q\circ\Phi(-)\)
\par\noindent\textbf{Corrigé:} \(\varinjlim_{X/\mathsf{Qis}^\ast} Q'\circ\Phi(-)\)
\par\noindent The right-derived limit formula likewise prints Q rather than Q′ before Φ. Its values lie in K(B), so the target-category quotient is Q′.
\smallskip
\item[\hypertarget{apparatus:FR-0306}{306}] Page \hyperlink{printpage:fr:305}{305}\quad\texttt{S45-SOL-1223}
\par\noindent\textbf{Imprimé:} \(\varinjlim_{\mathsf{Qis}^\ast/X}\)
\par\noindent\textbf{Corrigé:} \(\varprojlim_{\mathsf{Qis}^\ast/X}\)
\par\noindent The enlarged scan crop shows a right-pointing arrow beneath Lim in the left-derived formula. The dual construction is a limit over quasi-isomorphisms into X, hence an inverse limit; reverse the limit direction in corrected forms.
\smallskip
\item[\hypertarget{apparatus:FR-0307}{307}] Page \hyperlink{printpage:fr:305}{305}\quad\texttt{S45-SOL-1224}
\par\noindent\textbf{Imprimé:} \(Q\circ \Phi(-)\)
\par\noindent\textbf{Corrigé:} \(Q'\circ \Phi(-)\)
\par\noindent The left-derived formula also prints Q∘Φ. As in(c) and the right-derived formula, Q′ is required by the source and target categories.
\smallskip
\item[\hypertarget{apparatus:FR-0308}{308}] Page \hyperlink{printpage:fr:309}{309}\quad\texttt{S45-ASTRA-V-0040}
\par\noindent\textbf{Imprimé:} \(\eK(\cA)\)
\par\noindent\textbf{Corrigé:} \(\eK(\cA)^\circ\)
\par\noindent RHom(-,Y) is contravariant in X; the ensuing bifunctor has D(A)\textasciicircum{}circ explicitly. The scan omits the opposite-category superscript in this display. Accept the corrected-tree repair and propagate it to English.
\smallskip
\item[\hypertarget{apparatus:FR-0309}{309}] Page \hyperlink{printpage:fr:310}{310}\quad\texttt{S45-SOL-1235}
\par\noindent\textbf{Imprimé:} {\footnotesize\ttfamily\raggedright 0 \textbackslash{}\allowbreak{}ar@\{\allowbreak{}<-\allowbreak{}\}\allowbreak{}[\allowbreak{}r]\allowbreak{}}
\par\noindent\textbf{Corrigé:} {\footnotesize\ttfamily\raggedright 0 \textbackslash{}\allowbreak{}ar[\allowbreak{}r]\allowbreak{}}
\par\noindent The first arrow in the printed short exact sequence points left, from Y′ to0. Exactness in the stated order and the tensor sequence immediately below require0→Y′. Preserve the arrow reversal diplomatically.
\smallskip
\item[\hypertarget{apparatus:FR-0310}{310}] Page \hyperlink{printpage:fr:311}{311}\quad\texttt{S45-SOL-1238}
\par\noindent\textbf{Imprimé:} {\footnotesize\ttfamily\raggedright Soient \$X\$ un objet de \$\textbackslash{}\allowbreak{}eD\textasciicircum{}\allowbreak{}-\allowbreak{}(\allowbreak{}\textbackslash{}\allowbreak{}underline E)\allowbreak{}\$ et \$Y\$ un objet de  \$\textbackslash{}\allowbreak{}eD\textasciicircum{}\allowbreak{}+\allowbreak{}(\allowbreak{}\textbackslash{}\allowbreak{}underline\{\allowbreak{}E\}\allowbreak{}')\allowbreak{}\$; il existe un isomorphisme fonctoriel:  \textbackslash{}\allowbreak{}[\allowbreak{}\textbackslash{}\allowbreak{}xymatrix\{\allowbreak{}   \textbackslash{}\allowbreak{}hom\_\allowbreak{}\{\allowbreak{}\textbackslash{}\allowbreak{}eD(\allowbreak{}\textbackslash{}\allowbreak{}underline\{\allowbreak{}E\}\allowbreak{}')\allowbreak{}\}\allowbreak{}(\allowbreak{}\textbackslash{}\allowbreak{}eL\textasciicircum{}\allowbreak{}-\allowbreak{} f\textasciicircum{}\allowbreak{}\textbackslash{}\allowbreak{}ast(\allowbreak{}X)\allowbreak{},\allowbreak{}Y)\allowbreak{} \textbackslash{}\allowbreak{}ar[\allowbreak{}r]\allowbreak{}\textasciicircum{}\allowbreak{}-\allowbreak{}\textbackslash{}\allowbreak{}sim      \& \textbackslash{}\allowbreak{}hom\_\allowbreak{}\{\allowbreak{}\textbackslash{}\allowbreak{}eD(\allowbreak{}\textbackslash{}\allowbreak{}underline E)\allowbreak{}\}\allowbreak{}(\allowbreak{}X,\allowbreak{}\textbackslash{}\allowbreak{}eR\textasciicircum{}\allowbreak{}+\allowbreak{} f\_\allowbreak{}\textbackslash{}\allowbreak{}ast(\allowbreak{}Y)\allowbreak{})\allowbreak{} \textbackslash{}\allowbreak{}text\{\allowbreak{}.\}\allowbreak{} \}\allowbreak{}\textbackslash{}\allowbreak{}]\allowbreak{}}
\par\noindent\textbf{Corrigé:} {\footnotesize\ttfamily\raggedright Soient \$X\$ un objet de \$\textbackslash{}\allowbreak{}eD\textasciicircum{}\allowbreak{}-\allowbreak{}(\allowbreak{}\textbackslash{}\allowbreak{}underline\{\allowbreak{}E\}\allowbreak{}')\allowbreak{}\$ et \$Y\$ un objet de  \$\textbackslash{}\allowbreak{}eD\textasciicircum{}\allowbreak{}+\allowbreak{}(\allowbreak{}\textbackslash{}\allowbreak{}underline E)\allowbreak{}\$; il existe un isomorphisme fonctoriel:  \textbackslash{}\allowbreak{}[\allowbreak{}\textbackslash{}\allowbreak{}xymatrix\{\allowbreak{}   \textbackslash{}\allowbreak{}hom\_\allowbreak{}\{\allowbreak{}\textbackslash{}\allowbreak{}eD(\allowbreak{}\textbackslash{}\allowbreak{}underline E)\allowbreak{}\}\allowbreak{}(\allowbreak{}\textbackslash{}\allowbreak{}eL\textasciicircum{}\allowbreak{}-\allowbreak{} f\textasciicircum{}\allowbreak{}\textbackslash{}\allowbreak{}ast(\allowbreak{}X)\allowbreak{},\allowbreak{}Y)\allowbreak{} \textbackslash{}\allowbreak{}ar[\allowbreak{}r]\allowbreak{}\textasciicircum{}\allowbreak{}-\allowbreak{}\textbackslash{}\allowbreak{}sim      \& \textbackslash{}\allowbreak{}hom\_\allowbreak{}\{\allowbreak{}\textbackslash{}\allowbreak{}eD(\allowbreak{}\textbackslash{}\allowbreak{}underline\{\allowbreak{}E\}\allowbreak{}')\allowbreak{}\}\allowbreak{}(\allowbreak{}X,\allowbreak{}\textbackslash{}\allowbreak{}eR\textasciicircum{}\allowbreak{}+\allowbreak{} f\_\allowbreak{}\textbackslash{}\allowbreak{}ast(\allowbreak{}Y)\allowbreak{})\allowbreak{} \textbackslash{}\allowbreak{}text\{\allowbreak{}.\}\allowbreak{} \}\allowbreak{}\textbackslash{}\allowbreak{}]\allowbreak{}}
\par\noindent The scan reverses E and E′ in both object hypotheses and both Hom categories of the final adjunction. Since f*:E′→E and f\_*:E→E′, X belongs to D\textasciicircum{}-(E′), Y to D\textasciicircum{}+(E), the left Hom is in D(E), and the right Hom in D(E′). Correct all four positions consistently. The printed module-category notation is underlined E, restored in both branches.
\smallskip
\item[\hypertarget{apparatus:FR-0311}{311}] Page \hyperlink{printpage:fr:312}{312}\quad\texttt{S45-SOL-1241}
\par\noindent\textbf{Imprimé:} Arpès
\par\noindent\textbf{Corrigé:} Après
\par\noindent The printed word is Arpès, with the r and p transposed. Correct it to Après in French; English After is unchanged.
\smallskip
\end{description}
\endgroup
\fi
